% =====================================================================
% Chapter 4: The local-resonance gap
% Theory session, dispersion branch (Nov 2026)
% Requires: amsmath, amssymb, graphicx (for the figure)
% Include in the main file with  % =====================================================================
% Chapter 4: The local-resonance gap
% Theory session, dispersion branch (Nov 2026)
% Requires: amsmath, amssymb, graphicx (for the figure)
% Include in the main file with  % =====================================================================
% Chapter 4: The local-resonance gap
% Theory session, dispersion branch (Nov 2026)
% Requires: amsmath, amssymb, graphicx (for the figure)
% Include in the main file with  % =====================================================================
% Chapter 4: The local-resonance gap
% Theory session, dispersion branch (Nov 2026)
% Requires: amsmath, amssymb, graphicx (for the figure)
% Include in the main file with  \include{chapter_04_local_resonance}
% =====================================================================

\chapter{The local-resonance gap}
\label{ch:local_resonance}

This chapter replaces the rigid mass of Chapters~\ref{ch:two_level} and
\ref{ch:bragg_2d} by a real mass--spring resonator. The structure stays the
same: one attachment point per cell, coupling only through
$w(\mathbf{x}_0)$, and two-level models.

\begin{center}
\begin{tabular}{|l|p{0.45\textwidth}|l|}
\hline
Symbol & Meaning & Value / units \\
\hline
$m$ & resonator mass, one per cell & $98\,\mathrm{g}$ \\
$k_s$ & spring stiffness (subscript $s$ to avoid confusion with the wavevector $\mathbf{k}$) & $\approx 5.6\cdot 10^6\,\mathrm{N/m}$ \\
$\lambda_r = \omega_r^2 = k_s/m$ & squared angular tuning frequency & $f_r = 1200\,\mathrm{Hz}$ \\
$u$ & displacement of the resonator mass & m \\
$w_0 = w(\mathbf{x}_0)$ & plate deflection at the attachment point & m \\
$\mu = m/(\rho h A)$ & mass ratio & $0.25$ \\
$\lambda = \omega^2$ & \textbf{unknown}: $\omega^2$ of a mode of plate and resonator (one dot of the diagram) & $\mathrm{rad^2/s^2}$ \\
$\lambda_p(\mathbf{k})$ & $\omega^2$ of one bare plate mode at this $\mathbf{k}$ (the $\lambda_n^0$ of Chapter~\ref{ch:two_level}) & $\mathrm{rad^2/s^2}$ \\
$q$ & modal amplitude of that plate mode & $\sqrt{\mathrm{kg}}\cdot\mathrm{m}$ \\
\hline
\end{tabular}
\end{center}

% ---------------------------------------------------------------------
\section{The resonator as a frequency-dependent mass}

\paragraph{Equation of motion.}
The spring is stretched by $w_0 - u$, so $m\ddot{u} = k_s(w_0 - u)$. For
harmonic motion, $\ddot{u} = -\omega^2 u$:
\begin{align}
  -m\omega^2 u &= k_s w_0 - k_s u \\
  u\,(k_s - m\omega^2) &= k_s w_0 \\
  u &= \frac{k_s w_0}{k_s - m\omega^2} = w_0\,\frac{\omega_r^2}{\omega_r^2 - \omega^2},
\end{align}
where the last step divides numerator and denominator by $m$.

\paragraph{Force on the plate.}
The spring pulls the plate towards the resonator, $F = k_s(u - w_0)$. With
\begin{equation}
  u - w_0 = w_0\left[\frac{\omega_r^2}{\omega_r^2 - \omega^2} - 1\right]
          = w_0\,\frac{\omega^2}{\omega_r^2 - \omega^2}
\end{equation}
and $k_s = m\omega_r^2$ we get
\begin{equation}
  F = \omega^2 w_0 \cdot \frac{m\,\omega_r^2}{\omega_r^2 - \omega^2} .
\end{equation}
A rigid mass pushed on the plate with $F = m\omega^2 w_0$. The resonator gives
the same expression with $m$ replaced by the \emph{effective mass}
\begin{equation}
  \boxed{\;m_\mathrm{eff}(\omega) = \frac{m\,\omega_r^2}{\omega_r^2 - \omega^2}\;}
\end{equation}
Everything from the rigid-mass chapters carries over by replacing $m$ with
$m_\mathrm{eff}(\omega)$ and $\mu$ with
$\mu_\mathrm{eff}(\omega) = \mu\,\omega_r^2/(\omega_r^2 - \omega^2)$.

\begin{center}
\begin{tabular}{|l|l|p{0.22\textwidth}|p{0.3\textwidth}|}
\hline
Frequency & $m_\mathrm{eff}$ & Resonator motion & What the plate feels \\
\hline
$\omega \ll \omega_r$ & $\approx m$ & $u \approx w_0$, with the plate & a rigid mass \\
just below $\omega_r$ & $\to +\infty$ & $u \gg w_0$, in phase & a huge mass: the plate is pinned \\
just above $\omega_r$ & $\to -\infty$ & $u \gg w_0$, antiphase & the spring pulls the plate back against its motion \\
$\omega \gg \omega_r$ & $\approx -m\omega_r^2/\omega^2$ & $u \approx 0$, mass stands still & $F \to -k_s w_0$: a spring to a fixed point \\
\hline
\end{tabular}
\end{center}
A negative $m_\mathrm{eff}$ is not strange. It means that the force of the
resonator acts as a \emph{restoring} force (stiffness) and not as an inertial
one. The last row shows this most clearly.

% ---------------------------------------------------------------------
\section{Effective medium: smearing the resonators over the plate}

For waves much longer than $a$ only $\mathbf{G} = 0$ matters, and $m_\mathrm{eff}$
can be spread over the cell area: $\rho h \to \rho h + m_\mathrm{eff}/A$. The
Kirchhoff relation $D\kappa^4 = \rho h\,\omega^2$ becomes
\begin{equation}
  D\kappa^4 = \omega^2 \rho h \left[1 + \frac{\mu\,\omega_r^2}{\omega_r^2 - \omega^2}\right] .
\end{equation}
Bringing the bracket to a common denominator,
\begin{equation}
  1 + \frac{\mu\,\omega_r^2}{\omega_r^2 - \omega^2}
  = \frac{\omega_r^2 - \omega^2 + \mu\,\omega_r^2}{\omega_r^2 - \omega^2}
  = \frac{\omega_r^2(1 + \mu) - \omega^2}{\omega_r^2 - \omega^2},
\end{equation}
so that
\begin{equation}
  D\kappa^4 = \omega^2 \rho h\,\frac{\omega_r^2(1 + \mu) - \omega^2}{\omega_r^2 - \omega^2} .
  \label{eq:effective_medium}
\end{equation}
A real $\kappa$ needs a non-negative right-hand side:
\begin{center}
\begin{tabular}{|l|c|c|l|}
\hline
Range & Numerator & Denominator & $\kappa^4$ \\
\hline
$\omega < \omega_r$ & $+$ & $+$ & $+$: waves propagate \\
$\omega_r < \omega < \omega_r\sqrt{1+\mu}$ & $+$ & $-$ & $-$: no real $\kappa$, gap \\
$\omega > \omega_r\sqrt{1+\mu}$ & $-$ & $-$ & $+$: waves propagate \\
\hline
\end{tabular}
\end{center}
The effective-medium gap is $f_r \dots f_r\sqrt{1+\mu} = 1200 \dots
1342\,\mathrm{Hz}$. The FE result is $1031$--$1195\,\mathrm{Hz}$; the difference
is the topic of the next chapter.

% ---------------------------------------------------------------------
\section{The two-level model with a resonator}

\subsection{The modal amplitude $q$}
Keep one plate mode only, so the deflection is
\begin{equation}
  w(\mathbf{x}) = q\,\varphi(\mathbf{x}),
  \qquad
  \varphi(\mathbf{x}) = \frac{e^{i(\mathbf{k}+\mathbf{G})\cdot\mathbf{x}}}{\sqrt{\rho h A}} .
\end{equation}
$q$ is the \emph{modal amplitude}: how strongly this mode is present. Since
$\varphi$ has units $1/\sqrt{\mathrm{kg}}$, $q$ has units
$\sqrt{\mathrm{kg}}\cdot\mathrm{m}$. If $W$ is the deflection amplitude in
metres, $w = W e^{i(\mathbf{k}+\mathbf{G})\cdot\mathbf{x}}$ and
$q = W\sqrt{\rho h A}$. At the attachment point $w_0 = \varphi\, q$, where
$\varphi$ without an argument means $\varphi(\mathbf{x}_0)$. The unknowns are
collected in the vector $x = (q, u)^T$.

\subsection{Direct route: two force balances}

\paragraph{Row 2: Newton for the resonator.}
The spring stretch is $\varphi q - u$, so $-m\omega^2 u = k_s(\varphi q - u)$.
With $\lambda = \omega^2$:
\begin{equation}
  -k_s\varphi\, q + k_s u = \lambda\, m\, u .
  \label{eq:row_resonator}
\end{equation}

\paragraph{Row 1: the plate mode.}
This is the derivation of Chapter~\ref{ch:two_level} with another force. The
spring pulls on the plate at $\mathbf{x}_0$ with $F = k_s(u - w_0)$:
\begin{equation}
  \mathcal{L} w - \omega^2 \rho h\, w = F\,\delta(\mathbf{x} - \mathbf{x}_0) .
\end{equation}
Insert $w = q\varphi$, multiply by $\varphi^*$ and integrate over the cell. The
$\mathcal{L}$ term gives $\lambda_p q$, the $\rho h$ term $\lambda q$, and the
$\delta$ term $\varphi^* F = \varphi^* k_s(u - \varphi q)$. Hence
$\lambda_p q - \lambda q = k_s\varphi^* u - k_s|\varphi|^2 q$, or
\begin{equation}
  (\lambda_p + k_s|\varphi|^2)\, q - k_s\varphi^*\, u = \lambda\, q .
  \label{eq:row_plate}
\end{equation}

\paragraph{Matrix form.}
Eqs.~\eqref{eq:row_plate} and \eqref{eq:row_resonator} together read
$K x = \lambda M x$ with
\begin{equation}
  K = \begin{pmatrix} \lambda_p + k_s|\varphi|^2 & -k_s\varphi^* \\ -k_s\varphi & k_s \end{pmatrix},
  \qquad
  M = \begin{pmatrix} 1 & 0 \\ 0 & m \end{pmatrix} .
  \label{eq:KM_resonator}
\end{equation}
Each row of $K$ holds the coefficients of $q$ and $u$ on the left of one
equation; each row of $M$ the coefficients on the right. The $1$ in $M$ is the
modal mass of the plate mode, equal to one by the normalization. The entries of
$K$ have different units because $q$ and $u$ do.

\subsection{The same matrices from the energies}

\paragraph{Plate kinetic energy.} Kinetic energy is
$\tfrac12\cdot\text{mass}\cdot\text{velocity}^2$, summed over the plate, with
velocity amplitude $\omega w = \omega q\varphi$:
\begin{equation}
  T_\mathrm{plate} = \tfrac12 \int \rho h\, \omega^2 |q|^2 |\varphi|^2 \,\mathrm{d}A
  = \tfrac12\,\omega^2 |q|^2 \int \rho h\, |\varphi|^2\,\mathrm{d}A
  = \tfrac12\,\omega^2 |q|^2 .
\end{equation}

\paragraph{Resonator kinetic energy.} $T_\mathrm{res} = \tfrac12\, m\, (\omega u)^2 = \tfrac12\,\omega^2 m |u|^2$.

\paragraph{Plate strain energy.} A mode acts like a single oscillator of mass
one. From $\omega^2 = \text{stiffness}/\text{mass}$ its stiffness is
$\lambda_p$, so the energy is $\tfrac12\lambda_p|q|^2$. Formally,
$\tfrac12\int w^*\mathcal{L}w\,\mathrm{d}A = \tfrac12|q|^2\lambda_p\int\rho h|\varphi|^2\,\mathrm{d}A$.

\paragraph{Spring strain energy.} $\tfrac12\, k_s\cdot\text{stretch}^2 = \tfrac12\, k_s|u - \varphi q|^2$.

\paragraph{What multiplying out does.}
Any such energy can be written with a matrix, which is how $K$ and $M$ are
defined:
\begin{align}
  \text{strain energy} &= \tfrac12\, x^H K x = \tfrac12\left(K_{11}|q|^2 + K_{12}\,q^* u + K_{21}\,u^* q + K_{22}|u|^2\right), \\
  \text{kinetic energy} &= \tfrac12\,\omega^2\, x^H M x .
\end{align}
$K$ is read off by sorting the energy into the four products $|q|^2$, $q^*u$,
$u^*q$, $|u|^2$. Multiplying out does that sorting:
\begin{equation}
  |u - \varphi q|^2 = (u^* - \varphi^* q^*)(u - \varphi q)
  = |u|^2 - \varphi\, q\, u^* - \varphi^* q^* u + |\varphi|^2 |q|^2 .
\end{equation}
\begin{center}
\begin{tabular}{|l|l|l|}
\hline
Product & Coefficient in $2\times$ strain energy & Entry \\
\hline
$|q|^2$ & $\lambda_p + k_s|\varphi|^2$ & $K_{11}$ \\
$q^* u$ & $-k_s\varphi^*$ & $K_{12}$ \\
$u^* q$ & $-k_s\varphi$ & $K_{21}$ \\
$|u|^2$ & $k_s$ & $K_{22}$ \\
\hline
\end{tabular}
\end{center}
For $M$, $2\times$ kinetic energy$/\omega^2 = |q|^2 + m|u|^2$ gives
$M_{11} = 1$, $M_{22} = m$. $M$ has no off-diagonal entries because no kinetic
energy term contains both $q$ and $u$: the masses are connected only by the
spring, so the coupling sits in $K$. Energy sorting and force balance are two
routes to the same result; FE assembly uses the energy route.

\paragraph{Link to the code.}
The stretch is $r\,x$ with the row vector $r = (\varphi, -1)$, so the spring
part of $K$ is $k_s\, r^H r$. This is what the \texttt{augment\_*} functions
build: plate $K$ plus this rank-one term, and $m$ on the diagonal of the
resonator DOF. In the code the vector is real (FE basis weights), and the
phases enter through $T(\mathbf{k})$.

\subsection{The standard form $H$}
With the mass-normalized resonator coordinate $v = \sqrt{m}\,u$, $M$ becomes
the identity and $K$ becomes a Hermitian matrix $H$. For a plane wave,
$|\varphi|^2 = 1/(\rho h A)$, and with $\mathbf{x}_0 = 0$, $\varphi$ is real:
\begin{align}
  \text{plate diagonal:} \quad & k_s|\varphi|^2 = \frac{k_s}{\rho h A} = \frac{k_s}{m}\cdot\frac{m}{\rho h A} = \mu\lambda_r, \\
  \text{resonator diagonal:} \quad & \frac{k_s}{m} = \lambda_r, \\
  \text{off-diagonal:} \quad & \frac{k_s\varphi}{\sqrt{m}} = \frac{k_s}{\sqrt{m\,\rho h A}} = \frac{k_s}{m}\sqrt{\frac{m}{\rho h A}} = \sqrt{\mu}\,\lambda_r .
\end{align}
So
\begin{equation}
  \boxed{\;H = \begin{pmatrix} \lambda_p + \mu\lambda_r & -\sqrt{\mu}\,\lambda_r \\ -\sqrt{\mu}\,\lambda_r & \lambda_r \end{pmatrix}\;}
  \label{eq:H_resonator}
\end{equation}
\begin{itemize}
  \item $\lambda_p + \mu\lambda_r$: plate mode, slightly \emph{stiffened} by the
        attached spring;
  \item $\lambda_r$: the resonator alone on a fixed base;
  \item $c = -\sqrt{\mu}\,\lambda_r$: the coupling through the spring.
\end{itemize}
Unlike the rigid-mass matrix, $H$ does not contain the unknown $\lambda$. The
extra DOF makes the problem linear again, which is why the FE code augments
instead of using $m_\mathrm{eff}(\omega)$.

\subsection{Checks}

\paragraph{$\Gamma$, band 1 ($\lambda_p = 0$: the cell translates rigidly).}
$H = \lambda_r \begin{pmatrix} \mu & -\sqrt{\mu} \\ -\sqrt{\mu} & 1 \end{pmatrix}$
has trace $\lambda_r(1+\mu)$ and determinant $\lambda_r^2(\mu - \mu) = 0$. The
eigenvalues are $0$ (plate and resonator translate together) and
$\lambda_r(1+\mu)$, the upper edge of the effective-medium gap.

\paragraph{Short waves ($\lambda_p \gg \lambda_r$).}
Expanding as in Section~\ref{sec:far}, the lower branch tends to $\lambda_r$:
the plate is too stiff to move, and the resonator vibrates on a fixed base.

\paragraph{Back to the effective medium.}
$\det(H - \lambda I) = 0$ reads
\begin{equation}
  (\lambda_p + \mu\lambda_r - \lambda)(\lambda_r - \lambda) - \mu\lambda_r^2 = 0 .
\end{equation}
Expanding the product,
\begin{equation}
  \lambda_p(\lambda_r - \lambda) + \mu\lambda_r(\lambda_r - \lambda) - \lambda(\lambda_r - \lambda) - \mu\lambda_r^2 = 0 .
\end{equation}
The $\mu\lambda_r^2$ terms cancel:
\begin{equation}
  \lambda_p(\lambda_r - \lambda) - \mu\lambda_r\lambda - \lambda(\lambda_r - \lambda) = 0 .
\end{equation}
Solving for $\lambda_p$:
\begin{equation}
  \lambda_p = \lambda\,\frac{\lambda_r(1+\mu) - \lambda}{\lambda_r - \lambda} .
\end{equation}
With $\lambda_p = D\kappa^4/\rho h$ this is exactly
Eq.~\eqref{eq:effective_medium}. The two-level model with one plate mode
\emph{is} the effective medium.

\subsection{The avoided crossing}
The diagonal entries are equal where $\lambda_p + \mu\lambda_r = \lambda_r$,
i.e.\ $\lambda_p = \lambda_r(1-\mu)$. There the splitting is $\pm|c|$:
\begin{equation}
  \lambda = \lambda_r \pm \sqrt{\mu}\,\lambda_r = \lambda_r(1 \pm \sqrt{\mu}) .
\end{equation}
For $\mu = 0.25$: $f = f_r\sqrt{0.5} \approx 849\,\mathrm{Hz}$ and
$f = f_r\sqrt{1.5} \approx 1470\,\mathrm{Hz}$.

\paragraph{Splitting versus gap width.}
$\pm\sqrt{\mu}$ is the spacing of the two bands \emph{at that one
$\mathbf{k}$}. In the unfolded model the complete gap runs from the
\emph{maximum} of the lower band (tending to $f_r$) to the \emph{minimum} of the
upper band ($f_r\sqrt{1+\mu}$ at $\Gamma$). Splitting and gap width are
different things.

\paragraph{Why $\sqrt{\mu}$ here and $\mu$ for the rigid mass?}
In Chapter~\ref{ch:two_level} two \emph{plate} modes were coupled indirectly
through the inertia of the mass; both ends of the chain carried a factor
$m|\varphi|^2$-type term, so the coupling went as $\mu$. Here the resonator is
a degree of freedom of its own, attached directly by the spring. Each
normalized coordinate carries one square root of its mass, which gives
$\sqrt{\mu}$.

% ---------------------------------------------------------------------
\section{The flat band and the plate share}

\paragraph{Where the flat band comes from.}
The augmented problem has one more DOF per cell than the bare plate, so there
is one extra eigenvalue at every $\mathbf{k}$. Away from plate bands it is the
resonator vibrating while the plate moves little. The resonator is a
\emph{local} DOF, and $\mathbf{k}$ enters only through the phase factors of
the plate; if the plate moves little, the eigenvalue barely depends on
$\mathbf{k}$.

\paragraph{The slope of a band.}
$H$ in Eq.~\eqref{eq:H_resonator} depends on $\mathbf{k}$ only through
$\lambda_p$, in one entry. First-order perturbation theory then gives, for a
normalized eigenvector $(q, v)$,
\begin{equation}
  \frac{\mathrm{d}\lambda}{\mathrm{d}\lambda_p} = |q|^2 = \text{plate share of the mode},
\end{equation}
where \emph{plate share} means the fraction of the kinetic energy of the mode
that is in the plate. Therefore
\begin{equation}
  \text{slope of a band} = \text{plate share} \times \text{slope of the plate curve} .
\end{equation}
A band is flat where the plate share is small, or where the plate curve itself
is flat, and it bends where the plate share grows. This happens near the plate
curve, not simply at some distance from $\omega_r$.

\paragraph{Computing the plate share.}
Row 2 of Eq.~\eqref{eq:H_resonator} gives
$-\sqrt{\mu}\,\lambda_r\, q + \lambda_r v = \lambda v$, hence
\begin{equation}
  \frac{q}{v} = \frac{\lambda_r - \lambda}{\sqrt{\mu}\,\lambda_r},
  \qquad
  \text{plate share} = \frac{(q/v)^2}{1 + (q/v)^2} .
\end{equation}
At $X$ and $M$ the coupled combination has $|\varphi|^2 = 2/(\rho h A)$ and
$4/(\rho h A)$ (Chapter~\ref{ch:bragg_2d}), so in all formulas $\mu$ is
replaced by $2\mu$ and $4\mu$, and $\lambda_p$ by the bare value
($0.983\,\lambda_r$ at $X$, $3.93\,\lambda_r$ at $M$).

% ---------------------------------------------------------------------
\section{Reading the dispersion diagram}

\begin{figure}[htbp]
  \centering
  % \includegraphics[width=0.85\textwidth]{dispersion_vamm_1200.png}
  \caption{Dispersion diagram of the $0.1\,\mathrm{m}$ unit cell
    ($h = 5\,\mathrm{mm}$, steel) with one resonator per cell,
    $f_r = 1200\,\mathrm{Hz}$, $\mu = 0.25$, along
    $\Gamma$--$X$--$M$--$\Gamma$. Blue dots: FE Bloch solution. Solid grey:
    plane-wave expansion. Dashed grey: bare plate. Dotted red: resonator
    frequency $f_r$. The complete gap runs from the maximum of band 1 at $M$
    ($1031\,\mathrm{Hz}$) to the minimum of band 2 at $X$
    ($1195\,\mathrm{Hz}$).}
  \label{fig:dispersion_vamm_1200}
\end{figure}

The blue dots in Fig.~\ref{fig:dispersion_vamm_1200} are the actual modes. They
come from three ingredients:
\begin{enumerate}
  \item the grey dashed rising curve: the plate wave going out from $\Gamma$;
  \item the grey dashed falling curves: the folded plate waves coming down
        towards $X$ and $M$;
  \item the red dotted line: the resonator alone, at $1200\,\mathrm{Hz}$.
\end{enumerate}
Wherever these meet, the blue bands repel. Between meetings, each blue band
follows whichever ingredient is closest.

The cleanest place to see this is the right part of the plot, $M$--$\Gamma$.
The rising plate curve crosses $1200\,\mathrm{Hz}$ about $70\,\%$ of the way
from $\Gamma$ to $M$ (at $\kappa = \pi/a$); the diagonal entries of $H$ are equal
slightly earlier, at about two thirds of the way. Reading from $\Gamma$ towards
$M$:
\begin{itemize}
  \item \textbf{Band 1} starts on the grey curve, bends away below the crossing
        and flattens out.
  \item \textbf{Band 2} starts flat at about $1270\,\mathrm{Hz}$, stays flat,
        then bends upwards at the crossing and joins the grey curve.
\end{itemize}
The two bands swap roles at the crossing: band 1 starts plate-like and ends
resonator-like, band 2 does the opposite. This is the avoided crossing of the
previous section, and its model values $849$ and $1470\,\mathrm{Hz}$ fit the
plot at that position.

\paragraph{Why the bands are flat where they are.}
\begin{itemize}
  \item Flat at $1270\,\mathrm{Hz}$ near $\Gamma$: the plate share is low
        ($20\,\%$ in the model) and the plate curve itself is nearly flat
        there. Both reasons work together.
  \item Band 1 near $\Gamma$: the share is high ($80\,\%$), so the band follows
        the grey curve.
  \item The bending part of band 2 on $M$--$\Gamma$: the share grows towards
        $100\,\%$, and the band picks up the slope of the grey curve.
\end{itemize}

\paragraph{Plate share at the landmarks.}
The table extends the $2\times 2$ model to the coupled combinations at $X$ and
$M$; ``crossing'' is the avoided crossing on $M$--$\Gamma$. The shares are \emph{model estimates, not FE results}. The FE value
$2710\,\mathrm{Hz}$ was read off the plot.
\begin{center}
\begin{tabular}{|l|l|l|l|p{0.26\textwidth}|}
\hline
Point & Band & FE [Hz] & Model [Hz] & Plate share \\
\hline
$\Gamma$ & 1 & $0$ & $0$ & $80\,\%$ (resonators ride along) \\
$\Gamma$ & 2 (flat) & $\approx 1270$ & $1342$ & $20\,\%$ \\
crossing & 1 / 2 & -- & $849$ / $1470$ & $50\,\%$ / $50\,\%$ \\
$X$ & 1 & $836$ & $846$ & $34\,\%$ \\
$X$ & 2 & $1195$ & $1194$ & $100\,\%$ (node at the resonator) \\
$X$ & 3 & $1652$ & $1694$ & $66\,\%$ \\
$M$ & 1 & $1031$ & $1046$ & $5\,\%$ \\
$M$ & top of group & $\approx 2710$ & $2728$ & $95\,\%$ \\
\hline
\end{tabular}
\end{center}

\paragraph{Two points that are easy to get wrong.}
\begin{enumerate}
  \item Band 1 never reaches $f_r$ in this diagram, because the zone ends first.
        It tops out at $1031\,\mathrm{Hz}$ at $M$. The limit $f \to f_r$ exists
        only in the unfolded model; in the diagram it corresponds to band 1
        near $M$, with a plate share of $5\,\%$.
  \item The ``upper band at $\Gamma$'' of the two-level model \emph{is} the flat
        band of the plot. It sits at about $1270\,\mathrm{Hz}$, above the red
        line. Nothing in the diagram sits on $1200\,\mathrm{Hz}$.
\end{enumerate}

\paragraph{Outlook.}
The gap runs from the top of band 1 at $M$ to band 2 at $X$. That upper edge is
the $100\,\%$ plate mode, which does not feel the resonator at all. The next
chapter takes this up.

% ---------------------------------------------------------------------
\section{Check questions with answers}

\paragraph{Question 1.} \emph{Just above $f_r$, $m_\mathrm{eff}$ is large and
negative. What does the plate feel there, and why does that rule out a
propagating bending wave?}

\emph{Answer.} The resonator oscillates opposite to the plate and pulls it back,
so the usual propagating wave is stopped. Two refinements:
\begin{itemize}
  \item Nothing is dissipated, because the model is undamped. The wave becomes
        \emph{evanescent}: $\kappa^4 < 0$ makes $\kappa$ complex, so the plate
        still moves, but the motion decays from cell to cell. The energy is
        reflected, not absorbed.
  \item Negative $m_\mathrm{eff}$ alone is not enough. $m_\mathrm{eff}$ is
        negative at \emph{every} frequency above $f_r$, but the gap ends at
        $f_r\sqrt{1+\mu}$. What counts is the \emph{total} mass per area,
        $\rho h + m_\mathrm{eff}/A$. It is negative only while the negative
        mass of the resonator outweighs the mass of the plate, and it crosses
        zero exactly at $\omega^2 = \omega_r^2(1+\mu)$.
\end{itemize}

\paragraph{Question 2.} \emph{Why is the resonator band flat, and where would you
expect it not to be flat?}

\emph{Answer.} The plate moves little there, so the frequency barely depends on
$\mathbf{k}$. What decides flatness is not the distance from $\omega_r$ itself,
but how far the \emph{plate band at that $\mathbf{k}$} is from the resonator
level: slope of a band = plate share $\times$ slope of the plate curve. The
band bends at the avoided crossing, and it moves away from $\omega_r$
\emph{because} it hybridizes. The band near $\Gamma$ shows the difference: it
sits at a frequency different from $f_r$ ($1342\,\mathrm{Hz}$ in the model,
about $1270\,\mathrm{Hz}$ in FE) and is still nearly flat, because the plate
band is far below it there.

\paragraph{Question 3.} \emph{Along the lower band, the eigenvector goes from
mostly plate to mostly resonator. What does this mean for MAC tracking, and
which branch does a shaker on the finite plate see?}

\emph{Answer.} The 50/50 mix sits at the crossing, not near $f_r$; the mode
becomes resonator-dominated towards the end of band 1.

\emph{MAC tracking:} band 1 goes from $80\,\%$ plate at $\Gamma$ to $5\,\%$ at
$M$. That is a large change of character, but it happens gradually over many
$\mathbf{k}$-points, so the MAC between neighbouring points stays high and
tracking should follow the band without swapping labels. The MAC must use the
\emph{full vector including the resonator DOF}, with the augmented mass matrix;
on plate DOFs alone the comparison is somewhat weaker (lowest MAC of band 1:
$0.969$ instead of $0.990$), and leaving out the mass matrix is far worse ($0.619$).

\emph{Shaker:} the question has no clean one-branch answer. The shaker drives the
plate, so a mode is excited in proportion to its plate share. It sees band 1 well
below about $850\,\mathrm{Hz}$, band 1 near $M$ hardly at all, and the flat part
of band 2 only weakly (reasoning from the model, not measured). Near the gap
edges the upper branch keeps more plate motion than the lower one: at the upper
edge at $\Gamma$ the model gives a plate share of $\mu/(1+\mu) = 20\,\%$ and
$u/w = -1/\mu = -4$.

% =====================================================================

\chapter{The local-resonance gap}
\label{ch:local_resonance}

This chapter replaces the rigid mass of Chapters~\ref{ch:two_level} and
\ref{ch:bragg_2d} by a real mass--spring resonator. The structure stays the
same: one attachment point per cell, coupling only through
$w(\mathbf{x}_0)$, and two-level models.

\begin{center}
\begin{tabular}{|l|p{0.45\textwidth}|l|}
\hline
Symbol & Meaning & Value / units \\
\hline
$m$ & resonator mass, one per cell & $98\,\mathrm{g}$ \\
$k_s$ & spring stiffness (subscript $s$ to avoid confusion with the wavevector $\mathbf{k}$) & $\approx 5.6\cdot 10^6\,\mathrm{N/m}$ \\
$\lambda_r = \omega_r^2 = k_s/m$ & squared angular tuning frequency & $f_r = 1200\,\mathrm{Hz}$ \\
$u$ & displacement of the resonator mass & m \\
$w_0 = w(\mathbf{x}_0)$ & plate deflection at the attachment point & m \\
$\mu = m/(\rho h A)$ & mass ratio & $0.25$ \\
$\lambda = \omega^2$ & \textbf{unknown}: $\omega^2$ of a mode of plate and resonator (one dot of the diagram) & $\mathrm{rad^2/s^2}$ \\
$\lambda_p(\mathbf{k})$ & $\omega^2$ of one bare plate mode at this $\mathbf{k}$ (the $\lambda_n^0$ of Chapter~\ref{ch:two_level}) & $\mathrm{rad^2/s^2}$ \\
$q$ & modal amplitude of that plate mode & $\sqrt{\mathrm{kg}}\cdot\mathrm{m}$ \\
\hline
\end{tabular}
\end{center}

% ---------------------------------------------------------------------
\section{The resonator as a frequency-dependent mass}

\paragraph{Equation of motion.}
The spring is stretched by $w_0 - u$, so $m\ddot{u} = k_s(w_0 - u)$. For
harmonic motion, $\ddot{u} = -\omega^2 u$:
\begin{align}
  -m\omega^2 u &= k_s w_0 - k_s u \\
  u\,(k_s - m\omega^2) &= k_s w_0 \\
  u &= \frac{k_s w_0}{k_s - m\omega^2} = w_0\,\frac{\omega_r^2}{\omega_r^2 - \omega^2},
\end{align}
where the last step divides numerator and denominator by $m$.

\paragraph{Force on the plate.}
The spring pulls the plate towards the resonator, $F = k_s(u - w_0)$. With
\begin{equation}
  u - w_0 = w_0\left[\frac{\omega_r^2}{\omega_r^2 - \omega^2} - 1\right]
          = w_0\,\frac{\omega^2}{\omega_r^2 - \omega^2}
\end{equation}
and $k_s = m\omega_r^2$ we get
\begin{equation}
  F = \omega^2 w_0 \cdot \frac{m\,\omega_r^2}{\omega_r^2 - \omega^2} .
\end{equation}
A rigid mass pushed on the plate with $F = m\omega^2 w_0$. The resonator gives
the same expression with $m$ replaced by the \emph{effective mass}
\begin{equation}
  \boxed{\;m_\mathrm{eff}(\omega) = \frac{m\,\omega_r^2}{\omega_r^2 - \omega^2}\;}
\end{equation}
Everything from the rigid-mass chapters carries over by replacing $m$ with
$m_\mathrm{eff}(\omega)$ and $\mu$ with
$\mu_\mathrm{eff}(\omega) = \mu\,\omega_r^2/(\omega_r^2 - \omega^2)$.

\begin{center}
\begin{tabular}{|l|l|p{0.22\textwidth}|p{0.3\textwidth}|}
\hline
Frequency & $m_\mathrm{eff}$ & Resonator motion & What the plate feels \\
\hline
$\omega \ll \omega_r$ & $\approx m$ & $u \approx w_0$, with the plate & a rigid mass \\
just below $\omega_r$ & $\to +\infty$ & $u \gg w_0$, in phase & a huge mass: the plate is pinned \\
just above $\omega_r$ & $\to -\infty$ & $u \gg w_0$, antiphase & the spring pulls the plate back against its motion \\
$\omega \gg \omega_r$ & $\approx -m\omega_r^2/\omega^2$ & $u \approx 0$, mass stands still & $F \to -k_s w_0$: a spring to a fixed point \\
\hline
\end{tabular}
\end{center}
A negative $m_\mathrm{eff}$ is not strange. It means that the force of the
resonator acts as a \emph{restoring} force (stiffness) and not as an inertial
one. The last row shows this most clearly.

% ---------------------------------------------------------------------
\section{Effective medium: smearing the resonators over the plate}

For waves much longer than $a$ only $\mathbf{G} = 0$ matters, and $m_\mathrm{eff}$
can be spread over the cell area: $\rho h \to \rho h + m_\mathrm{eff}/A$. The
Kirchhoff relation $D\kappa^4 = \rho h\,\omega^2$ becomes
\begin{equation}
  D\kappa^4 = \omega^2 \rho h \left[1 + \frac{\mu\,\omega_r^2}{\omega_r^2 - \omega^2}\right] .
\end{equation}
Bringing the bracket to a common denominator,
\begin{equation}
  1 + \frac{\mu\,\omega_r^2}{\omega_r^2 - \omega^2}
  = \frac{\omega_r^2 - \omega^2 + \mu\,\omega_r^2}{\omega_r^2 - \omega^2}
  = \frac{\omega_r^2(1 + \mu) - \omega^2}{\omega_r^2 - \omega^2},
\end{equation}
so that
\begin{equation}
  D\kappa^4 = \omega^2 \rho h\,\frac{\omega_r^2(1 + \mu) - \omega^2}{\omega_r^2 - \omega^2} .
  \label{eq:effective_medium}
\end{equation}
A real $\kappa$ needs a non-negative right-hand side:
\begin{center}
\begin{tabular}{|l|c|c|l|}
\hline
Range & Numerator & Denominator & $\kappa^4$ \\
\hline
$\omega < \omega_r$ & $+$ & $+$ & $+$: waves propagate \\
$\omega_r < \omega < \omega_r\sqrt{1+\mu}$ & $+$ & $-$ & $-$: no real $\kappa$, gap \\
$\omega > \omega_r\sqrt{1+\mu}$ & $-$ & $-$ & $+$: waves propagate \\
\hline
\end{tabular}
\end{center}
The effective-medium gap is $f_r \dots f_r\sqrt{1+\mu} = 1200 \dots
1342\,\mathrm{Hz}$. The FE result is $1031$--$1195\,\mathrm{Hz}$; the difference
is the topic of the next chapter.

% ---------------------------------------------------------------------
\section{The two-level model with a resonator}

\subsection{The modal amplitude $q$}
Keep one plate mode only, so the deflection is
\begin{equation}
  w(\mathbf{x}) = q\,\varphi(\mathbf{x}),
  \qquad
  \varphi(\mathbf{x}) = \frac{e^{i(\mathbf{k}+\mathbf{G})\cdot\mathbf{x}}}{\sqrt{\rho h A}} .
\end{equation}
$q$ is the \emph{modal amplitude}: how strongly this mode is present. Since
$\varphi$ has units $1/\sqrt{\mathrm{kg}}$, $q$ has units
$\sqrt{\mathrm{kg}}\cdot\mathrm{m}$. If $W$ is the deflection amplitude in
metres, $w = W e^{i(\mathbf{k}+\mathbf{G})\cdot\mathbf{x}}$ and
$q = W\sqrt{\rho h A}$. At the attachment point $w_0 = \varphi\, q$, where
$\varphi$ without an argument means $\varphi(\mathbf{x}_0)$. The unknowns are
collected in the vector $x = (q, u)^T$.

\subsection{Direct route: two force balances}

\paragraph{Row 2: Newton for the resonator.}
The spring stretch is $\varphi q - u$, so $-m\omega^2 u = k_s(\varphi q - u)$.
With $\lambda = \omega^2$:
\begin{equation}
  -k_s\varphi\, q + k_s u = \lambda\, m\, u .
  \label{eq:row_resonator}
\end{equation}

\paragraph{Row 1: the plate mode.}
This is the derivation of Chapter~\ref{ch:two_level} with another force. The
spring pulls on the plate at $\mathbf{x}_0$ with $F = k_s(u - w_0)$:
\begin{equation}
  \mathcal{L} w - \omega^2 \rho h\, w = F\,\delta(\mathbf{x} - \mathbf{x}_0) .
\end{equation}
Insert $w = q\varphi$, multiply by $\varphi^*$ and integrate over the cell. The
$\mathcal{L}$ term gives $\lambda_p q$, the $\rho h$ term $\lambda q$, and the
$\delta$ term $\varphi^* F = \varphi^* k_s(u - \varphi q)$. Hence
$\lambda_p q - \lambda q = k_s\varphi^* u - k_s|\varphi|^2 q$, or
\begin{equation}
  (\lambda_p + k_s|\varphi|^2)\, q - k_s\varphi^*\, u = \lambda\, q .
  \label{eq:row_plate}
\end{equation}

\paragraph{Matrix form.}
Eqs.~\eqref{eq:row_plate} and \eqref{eq:row_resonator} together read
$K x = \lambda M x$ with
\begin{equation}
  K = \begin{pmatrix} \lambda_p + k_s|\varphi|^2 & -k_s\varphi^* \\ -k_s\varphi & k_s \end{pmatrix},
  \qquad
  M = \begin{pmatrix} 1 & 0 \\ 0 & m \end{pmatrix} .
  \label{eq:KM_resonator}
\end{equation}
Each row of $K$ holds the coefficients of $q$ and $u$ on the left of one
equation; each row of $M$ the coefficients on the right. The $1$ in $M$ is the
modal mass of the plate mode, equal to one by the normalization. The entries of
$K$ have different units because $q$ and $u$ do.

\subsection{The same matrices from the energies}

\paragraph{Plate kinetic energy.} Kinetic energy is
$\tfrac12\cdot\text{mass}\cdot\text{velocity}^2$, summed over the plate, with
velocity amplitude $\omega w = \omega q\varphi$:
\begin{equation}
  T_\mathrm{plate} = \tfrac12 \int \rho h\, \omega^2 |q|^2 |\varphi|^2 \,\mathrm{d}A
  = \tfrac12\,\omega^2 |q|^2 \int \rho h\, |\varphi|^2\,\mathrm{d}A
  = \tfrac12\,\omega^2 |q|^2 .
\end{equation}

\paragraph{Resonator kinetic energy.} $T_\mathrm{res} = \tfrac12\, m\, (\omega u)^2 = \tfrac12\,\omega^2 m |u|^2$.

\paragraph{Plate strain energy.} A mode acts like a single oscillator of mass
one. From $\omega^2 = \text{stiffness}/\text{mass}$ its stiffness is
$\lambda_p$, so the energy is $\tfrac12\lambda_p|q|^2$. Formally,
$\tfrac12\int w^*\mathcal{L}w\,\mathrm{d}A = \tfrac12|q|^2\lambda_p\int\rho h|\varphi|^2\,\mathrm{d}A$.

\paragraph{Spring strain energy.} $\tfrac12\, k_s\cdot\text{stretch}^2 = \tfrac12\, k_s|u - \varphi q|^2$.

\paragraph{What multiplying out does.}
Any such energy can be written with a matrix, which is how $K$ and $M$ are
defined:
\begin{align}
  \text{strain energy} &= \tfrac12\, x^H K x = \tfrac12\left(K_{11}|q|^2 + K_{12}\,q^* u + K_{21}\,u^* q + K_{22}|u|^2\right), \\
  \text{kinetic energy} &= \tfrac12\,\omega^2\, x^H M x .
\end{align}
$K$ is read off by sorting the energy into the four products $|q|^2$, $q^*u$,
$u^*q$, $|u|^2$. Multiplying out does that sorting:
\begin{equation}
  |u - \varphi q|^2 = (u^* - \varphi^* q^*)(u - \varphi q)
  = |u|^2 - \varphi\, q\, u^* - \varphi^* q^* u + |\varphi|^2 |q|^2 .
\end{equation}
\begin{center}
\begin{tabular}{|l|l|l|}
\hline
Product & Coefficient in $2\times$ strain energy & Entry \\
\hline
$|q|^2$ & $\lambda_p + k_s|\varphi|^2$ & $K_{11}$ \\
$q^* u$ & $-k_s\varphi^*$ & $K_{12}$ \\
$u^* q$ & $-k_s\varphi$ & $K_{21}$ \\
$|u|^2$ & $k_s$ & $K_{22}$ \\
\hline
\end{tabular}
\end{center}
For $M$, $2\times$ kinetic energy$/\omega^2 = |q|^2 + m|u|^2$ gives
$M_{11} = 1$, $M_{22} = m$. $M$ has no off-diagonal entries because no kinetic
energy term contains both $q$ and $u$: the masses are connected only by the
spring, so the coupling sits in $K$. Energy sorting and force balance are two
routes to the same result; FE assembly uses the energy route.

\paragraph{Link to the code.}
The stretch is $r\,x$ with the row vector $r = (\varphi, -1)$, so the spring
part of $K$ is $k_s\, r^H r$. This is what the \texttt{augment\_*} functions
build: plate $K$ plus this rank-one term, and $m$ on the diagonal of the
resonator DOF. In the code the vector is real (FE basis weights), and the
phases enter through $T(\mathbf{k})$.

\subsection{The standard form $H$}
With the mass-normalized resonator coordinate $v = \sqrt{m}\,u$, $M$ becomes
the identity and $K$ becomes a Hermitian matrix $H$. For a plane wave,
$|\varphi|^2 = 1/(\rho h A)$, and with $\mathbf{x}_0 = 0$, $\varphi$ is real:
\begin{align}
  \text{plate diagonal:} \quad & k_s|\varphi|^2 = \frac{k_s}{\rho h A} = \frac{k_s}{m}\cdot\frac{m}{\rho h A} = \mu\lambda_r, \\
  \text{resonator diagonal:} \quad & \frac{k_s}{m} = \lambda_r, \\
  \text{off-diagonal:} \quad & \frac{k_s\varphi}{\sqrt{m}} = \frac{k_s}{\sqrt{m\,\rho h A}} = \frac{k_s}{m}\sqrt{\frac{m}{\rho h A}} = \sqrt{\mu}\,\lambda_r .
\end{align}
So
\begin{equation}
  \boxed{\;H = \begin{pmatrix} \lambda_p + \mu\lambda_r & -\sqrt{\mu}\,\lambda_r \\ -\sqrt{\mu}\,\lambda_r & \lambda_r \end{pmatrix}\;}
  \label{eq:H_resonator}
\end{equation}
\begin{itemize}
  \item $\lambda_p + \mu\lambda_r$: plate mode, slightly \emph{stiffened} by the
        attached spring;
  \item $\lambda_r$: the resonator alone on a fixed base;
  \item $c = -\sqrt{\mu}\,\lambda_r$: the coupling through the spring.
\end{itemize}
Unlike the rigid-mass matrix, $H$ does not contain the unknown $\lambda$. The
extra DOF makes the problem linear again, which is why the FE code augments
instead of using $m_\mathrm{eff}(\omega)$.

\subsection{Checks}

\paragraph{$\Gamma$, band 1 ($\lambda_p = 0$: the cell translates rigidly).}
$H = \lambda_r \begin{pmatrix} \mu & -\sqrt{\mu} \\ -\sqrt{\mu} & 1 \end{pmatrix}$
has trace $\lambda_r(1+\mu)$ and determinant $\lambda_r^2(\mu - \mu) = 0$. The
eigenvalues are $0$ (plate and resonator translate together) and
$\lambda_r(1+\mu)$, the upper edge of the effective-medium gap.

\paragraph{Short waves ($\lambda_p \gg \lambda_r$).}
Expanding as in Section~\ref{sec:far}, the lower branch tends to $\lambda_r$:
the plate is too stiff to move, and the resonator vibrates on a fixed base.

\paragraph{Back to the effective medium.}
$\det(H - \lambda I) = 0$ reads
\begin{equation}
  (\lambda_p + \mu\lambda_r - \lambda)(\lambda_r - \lambda) - \mu\lambda_r^2 = 0 .
\end{equation}
Expanding the product,
\begin{equation}
  \lambda_p(\lambda_r - \lambda) + \mu\lambda_r(\lambda_r - \lambda) - \lambda(\lambda_r - \lambda) - \mu\lambda_r^2 = 0 .
\end{equation}
The $\mu\lambda_r^2$ terms cancel:
\begin{equation}
  \lambda_p(\lambda_r - \lambda) - \mu\lambda_r\lambda - \lambda(\lambda_r - \lambda) = 0 .
\end{equation}
Solving for $\lambda_p$:
\begin{equation}
  \lambda_p = \lambda\,\frac{\lambda_r(1+\mu) - \lambda}{\lambda_r - \lambda} .
\end{equation}
With $\lambda_p = D\kappa^4/\rho h$ this is exactly
Eq.~\eqref{eq:effective_medium}. The two-level model with one plate mode
\emph{is} the effective medium.

\subsection{The avoided crossing}
The diagonal entries are equal where $\lambda_p + \mu\lambda_r = \lambda_r$,
i.e.\ $\lambda_p = \lambda_r(1-\mu)$. There the splitting is $\pm|c|$:
\begin{equation}
  \lambda = \lambda_r \pm \sqrt{\mu}\,\lambda_r = \lambda_r(1 \pm \sqrt{\mu}) .
\end{equation}
For $\mu = 0.25$: $f = f_r\sqrt{0.5} \approx 849\,\mathrm{Hz}$ and
$f = f_r\sqrt{1.5} \approx 1470\,\mathrm{Hz}$.

\paragraph{Splitting versus gap width.}
$\pm\sqrt{\mu}$ is the spacing of the two bands \emph{at that one
$\mathbf{k}$}. In the unfolded model the complete gap runs from the
\emph{maximum} of the lower band (tending to $f_r$) to the \emph{minimum} of the
upper band ($f_r\sqrt{1+\mu}$ at $\Gamma$). Splitting and gap width are
different things.

\paragraph{Why $\sqrt{\mu}$ here and $\mu$ for the rigid mass?}
In Chapter~\ref{ch:two_level} two \emph{plate} modes were coupled indirectly
through the inertia of the mass; both ends of the chain carried a factor
$m|\varphi|^2$-type term, so the coupling went as $\mu$. Here the resonator is
a degree of freedom of its own, attached directly by the spring. Each
normalized coordinate carries one square root of its mass, which gives
$\sqrt{\mu}$.

% ---------------------------------------------------------------------
\section{The flat band and the plate share}

\paragraph{Where the flat band comes from.}
The augmented problem has one more DOF per cell than the bare plate, so there
is one extra eigenvalue at every $\mathbf{k}$. Away from plate bands it is the
resonator vibrating while the plate moves little. The resonator is a
\emph{local} DOF, and $\mathbf{k}$ enters only through the phase factors of
the plate; if the plate moves little, the eigenvalue barely depends on
$\mathbf{k}$.

\paragraph{The slope of a band.}
$H$ in Eq.~\eqref{eq:H_resonator} depends on $\mathbf{k}$ only through
$\lambda_p$, in one entry. First-order perturbation theory then gives, for a
normalized eigenvector $(q, v)$,
\begin{equation}
  \frac{\mathrm{d}\lambda}{\mathrm{d}\lambda_p} = |q|^2 = \text{plate share of the mode},
\end{equation}
where \emph{plate share} means the fraction of the kinetic energy of the mode
that is in the plate. Therefore
\begin{equation}
  \text{slope of a band} = \text{plate share} \times \text{slope of the plate curve} .
\end{equation}
A band is flat where the plate share is small, or where the plate curve itself
is flat, and it bends where the plate share grows. This happens near the plate
curve, not simply at some distance from $\omega_r$.

\paragraph{Computing the plate share.}
Row 2 of Eq.~\eqref{eq:H_resonator} gives
$-\sqrt{\mu}\,\lambda_r\, q + \lambda_r v = \lambda v$, hence
\begin{equation}
  \frac{q}{v} = \frac{\lambda_r - \lambda}{\sqrt{\mu}\,\lambda_r},
  \qquad
  \text{plate share} = \frac{(q/v)^2}{1 + (q/v)^2} .
\end{equation}
At $X$ and $M$ the coupled combination has $|\varphi|^2 = 2/(\rho h A)$ and
$4/(\rho h A)$ (Chapter~\ref{ch:bragg_2d}), so in all formulas $\mu$ is
replaced by $2\mu$ and $4\mu$, and $\lambda_p$ by the bare value
($0.983\,\lambda_r$ at $X$, $3.93\,\lambda_r$ at $M$).

% ---------------------------------------------------------------------
\section{Reading the dispersion diagram}

\begin{figure}[htbp]
  \centering
  % \includegraphics[width=0.85\textwidth]{dispersion_vamm_1200.png}
  \caption{Dispersion diagram of the $0.1\,\mathrm{m}$ unit cell
    ($h = 5\,\mathrm{mm}$, steel) with one resonator per cell,
    $f_r = 1200\,\mathrm{Hz}$, $\mu = 0.25$, along
    $\Gamma$--$X$--$M$--$\Gamma$. Blue dots: FE Bloch solution. Solid grey:
    plane-wave expansion. Dashed grey: bare plate. Dotted red: resonator
    frequency $f_r$. The complete gap runs from the maximum of band 1 at $M$
    ($1031\,\mathrm{Hz}$) to the minimum of band 2 at $X$
    ($1195\,\mathrm{Hz}$).}
  \label{fig:dispersion_vamm_1200}
\end{figure}

The blue dots in Fig.~\ref{fig:dispersion_vamm_1200} are the actual modes. They
come from three ingredients:
\begin{enumerate}
  \item the grey dashed rising curve: the plate wave going out from $\Gamma$;
  \item the grey dashed falling curves: the folded plate waves coming down
        towards $X$ and $M$;
  \item the red dotted line: the resonator alone, at $1200\,\mathrm{Hz}$.
\end{enumerate}
Wherever these meet, the blue bands repel. Between meetings, each blue band
follows whichever ingredient is closest.

The cleanest place to see this is the right part of the plot, $M$--$\Gamma$.
The rising plate curve crosses $1200\,\mathrm{Hz}$ about $70\,\%$ of the way
from $\Gamma$ to $M$ (at $\kappa = \pi/a$); the diagonal entries of $H$ are equal
slightly earlier, at about two thirds of the way. Reading from $\Gamma$ towards
$M$:
\begin{itemize}
  \item \textbf{Band 1} starts on the grey curve, bends away below the crossing
        and flattens out.
  \item \textbf{Band 2} starts flat at about $1270\,\mathrm{Hz}$, stays flat,
        then bends upwards at the crossing and joins the grey curve.
\end{itemize}
The two bands swap roles at the crossing: band 1 starts plate-like and ends
resonator-like, band 2 does the opposite. This is the avoided crossing of the
previous section, and its model values $849$ and $1470\,\mathrm{Hz}$ fit the
plot at that position.

\paragraph{Why the bands are flat where they are.}
\begin{itemize}
  \item Flat at $1270\,\mathrm{Hz}$ near $\Gamma$: the plate share is low
        ($20\,\%$ in the model) and the plate curve itself is nearly flat
        there. Both reasons work together.
  \item Band 1 near $\Gamma$: the share is high ($80\,\%$), so the band follows
        the grey curve.
  \item The bending part of band 2 on $M$--$\Gamma$: the share grows towards
        $100\,\%$, and the band picks up the slope of the grey curve.
\end{itemize}

\paragraph{Plate share at the landmarks.}
The table extends the $2\times 2$ model to the coupled combinations at $X$ and
$M$; ``crossing'' is the avoided crossing on $M$--$\Gamma$. The shares are \emph{model estimates, not FE results}. The FE value
$2710\,\mathrm{Hz}$ was read off the plot.
\begin{center}
\begin{tabular}{|l|l|l|l|p{0.26\textwidth}|}
\hline
Point & Band & FE [Hz] & Model [Hz] & Plate share \\
\hline
$\Gamma$ & 1 & $0$ & $0$ & $80\,\%$ (resonators ride along) \\
$\Gamma$ & 2 (flat) & $\approx 1270$ & $1342$ & $20\,\%$ \\
crossing & 1 / 2 & -- & $849$ / $1470$ & $50\,\%$ / $50\,\%$ \\
$X$ & 1 & $836$ & $846$ & $34\,\%$ \\
$X$ & 2 & $1195$ & $1194$ & $100\,\%$ (node at the resonator) \\
$X$ & 3 & $1652$ & $1694$ & $66\,\%$ \\
$M$ & 1 & $1031$ & $1046$ & $5\,\%$ \\
$M$ & top of group & $\approx 2710$ & $2728$ & $95\,\%$ \\
\hline
\end{tabular}
\end{center}

\paragraph{Two points that are easy to get wrong.}
\begin{enumerate}
  \item Band 1 never reaches $f_r$ in this diagram, because the zone ends first.
        It tops out at $1031\,\mathrm{Hz}$ at $M$. The limit $f \to f_r$ exists
        only in the unfolded model; in the diagram it corresponds to band 1
        near $M$, with a plate share of $5\,\%$.
  \item The ``upper band at $\Gamma$'' of the two-level model \emph{is} the flat
        band of the plot. It sits at about $1270\,\mathrm{Hz}$, above the red
        line. Nothing in the diagram sits on $1200\,\mathrm{Hz}$.
\end{enumerate}

\paragraph{Outlook.}
The gap runs from the top of band 1 at $M$ to band 2 at $X$. That upper edge is
the $100\,\%$ plate mode, which does not feel the resonator at all. The next
chapter takes this up.

% ---------------------------------------------------------------------
\section{Check questions with answers}

\paragraph{Question 1.} \emph{Just above $f_r$, $m_\mathrm{eff}$ is large and
negative. What does the plate feel there, and why does that rule out a
propagating bending wave?}

\emph{Answer.} The resonator oscillates opposite to the plate and pulls it back,
so the usual propagating wave is stopped. Two refinements:
\begin{itemize}
  \item Nothing is dissipated, because the model is undamped. The wave becomes
        \emph{evanescent}: $\kappa^4 < 0$ makes $\kappa$ complex, so the plate
        still moves, but the motion decays from cell to cell. The energy is
        reflected, not absorbed.
  \item Negative $m_\mathrm{eff}$ alone is not enough. $m_\mathrm{eff}$ is
        negative at \emph{every} frequency above $f_r$, but the gap ends at
        $f_r\sqrt{1+\mu}$. What counts is the \emph{total} mass per area,
        $\rho h + m_\mathrm{eff}/A$. It is negative only while the negative
        mass of the resonator outweighs the mass of the plate, and it crosses
        zero exactly at $\omega^2 = \omega_r^2(1+\mu)$.
\end{itemize}

\paragraph{Question 2.} \emph{Why is the resonator band flat, and where would you
expect it not to be flat?}

\emph{Answer.} The plate moves little there, so the frequency barely depends on
$\mathbf{k}$. What decides flatness is not the distance from $\omega_r$ itself,
but how far the \emph{plate band at that $\mathbf{k}$} is from the resonator
level: slope of a band = plate share $\times$ slope of the plate curve. The
band bends at the avoided crossing, and it moves away from $\omega_r$
\emph{because} it hybridizes. The band near $\Gamma$ shows the difference: it
sits at a frequency different from $f_r$ ($1342\,\mathrm{Hz}$ in the model,
about $1270\,\mathrm{Hz}$ in FE) and is still nearly flat, because the plate
band is far below it there.

\paragraph{Question 3.} \emph{Along the lower band, the eigenvector goes from
mostly plate to mostly resonator. What does this mean for MAC tracking, and
which branch does a shaker on the finite plate see?}

\emph{Answer.} The 50/50 mix sits at the crossing, not near $f_r$; the mode
becomes resonator-dominated towards the end of band 1.

\emph{MAC tracking:} band 1 goes from $80\,\%$ plate at $\Gamma$ to $5\,\%$ at
$M$. That is a large change of character, but it happens gradually over many
$\mathbf{k}$-points, so the MAC between neighbouring points stays high and
tracking should follow the band without swapping labels. The MAC must use the
\emph{full vector including the resonator DOF}, with the augmented mass matrix;
on plate DOFs alone the comparison is somewhat weaker (lowest MAC of band 1:
$0.969$ instead of $0.990$), and leaving out the mass matrix is far worse ($0.619$).

\emph{Shaker:} the question has no clean one-branch answer. The shaker drives the
plate, so a mode is excited in proportion to its plate share. It sees band 1 well
below about $850\,\mathrm{Hz}$, band 1 near $M$ hardly at all, and the flat part
of band 2 only weakly (reasoning from the model, not measured). Near the gap
edges the upper branch keeps more plate motion than the lower one: at the upper
edge at $\Gamma$ the model gives a plate share of $\mu/(1+\mu) = 20\,\%$ and
$u/w = -1/\mu = -4$.

% =====================================================================

\chapter{The local-resonance gap}
\label{ch:local_resonance}

This chapter replaces the rigid mass of Chapters~\ref{ch:two_level} and
\ref{ch:bragg_2d} by a real mass--spring resonator. The structure stays the
same: one attachment point per cell, coupling only through
$w(\mathbf{x}_0)$, and two-level models.

\begin{center}
\begin{tabular}{|l|p{0.45\textwidth}|l|}
\hline
Symbol & Meaning & Value / units \\
\hline
$m$ & resonator mass, one per cell & $98\,\mathrm{g}$ \\
$k_s$ & spring stiffness (subscript $s$ to avoid confusion with the wavevector $\mathbf{k}$) & $\approx 5.6\cdot 10^6\,\mathrm{N/m}$ \\
$\lambda_r = \omega_r^2 = k_s/m$ & squared angular tuning frequency & $f_r = 1200\,\mathrm{Hz}$ \\
$u$ & displacement of the resonator mass & m \\
$w_0 = w(\mathbf{x}_0)$ & plate deflection at the attachment point & m \\
$\mu = m/(\rho h A)$ & mass ratio & $0.25$ \\
$\lambda = \omega^2$ & \textbf{unknown}: $\omega^2$ of a mode of plate and resonator (one dot of the diagram) & $\mathrm{rad^2/s^2}$ \\
$\lambda_p(\mathbf{k})$ & $\omega^2$ of one bare plate mode at this $\mathbf{k}$ (the $\lambda_n^0$ of Chapter~\ref{ch:two_level}) & $\mathrm{rad^2/s^2}$ \\
$q$ & modal amplitude of that plate mode & $\sqrt{\mathrm{kg}}\cdot\mathrm{m}$ \\
\hline
\end{tabular}
\end{center}

% ---------------------------------------------------------------------
\section{The resonator as a frequency-dependent mass}

\paragraph{Equation of motion.}
The spring is stretched by $w_0 - u$, so $m\ddot{u} = k_s(w_0 - u)$. For
harmonic motion, $\ddot{u} = -\omega^2 u$:
\begin{align}
  -m\omega^2 u &= k_s w_0 - k_s u \\
  u\,(k_s - m\omega^2) &= k_s w_0 \\
  u &= \frac{k_s w_0}{k_s - m\omega^2} = w_0\,\frac{\omega_r^2}{\omega_r^2 - \omega^2},
\end{align}
where the last step divides numerator and denominator by $m$.

\paragraph{Force on the plate.}
The spring pulls the plate towards the resonator, $F = k_s(u - w_0)$. With
\begin{equation}
  u - w_0 = w_0\left[\frac{\omega_r^2}{\omega_r^2 - \omega^2} - 1\right]
          = w_0\,\frac{\omega^2}{\omega_r^2 - \omega^2}
\end{equation}
and $k_s = m\omega_r^2$ we get
\begin{equation}
  F = \omega^2 w_0 \cdot \frac{m\,\omega_r^2}{\omega_r^2 - \omega^2} .
\end{equation}
A rigid mass pushed on the plate with $F = m\omega^2 w_0$. The resonator gives
the same expression with $m$ replaced by the \emph{effective mass}
\begin{equation}
  \boxed{\;m_\mathrm{eff}(\omega) = \frac{m\,\omega_r^2}{\omega_r^2 - \omega^2}\;}
\end{equation}
Everything from the rigid-mass chapters carries over by replacing $m$ with
$m_\mathrm{eff}(\omega)$ and $\mu$ with
$\mu_\mathrm{eff}(\omega) = \mu\,\omega_r^2/(\omega_r^2 - \omega^2)$.

\begin{center}
\begin{tabular}{|l|l|p{0.22\textwidth}|p{0.3\textwidth}|}
\hline
Frequency & $m_\mathrm{eff}$ & Resonator motion & What the plate feels \\
\hline
$\omega \ll \omega_r$ & $\approx m$ & $u \approx w_0$, with the plate & a rigid mass \\
just below $\omega_r$ & $\to +\infty$ & $u \gg w_0$, in phase & a huge mass: the plate is pinned \\
just above $\omega_r$ & $\to -\infty$ & $u \gg w_0$, antiphase & the spring pulls the plate back against its motion \\
$\omega \gg \omega_r$ & $\approx -m\omega_r^2/\omega^2$ & $u \approx 0$, mass stands still & $F \to -k_s w_0$: a spring to a fixed point \\
\hline
\end{tabular}
\end{center}
A negative $m_\mathrm{eff}$ is not strange. It means that the force of the
resonator acts as a \emph{restoring} force (stiffness) and not as an inertial
one. The last row shows this most clearly.

% ---------------------------------------------------------------------
\section{Effective medium: smearing the resonators over the plate}

For waves much longer than $a$ only $\mathbf{G} = 0$ matters, and $m_\mathrm{eff}$
can be spread over the cell area: $\rho h \to \rho h + m_\mathrm{eff}/A$. The
Kirchhoff relation $D\kappa^4 = \rho h\,\omega^2$ becomes
\begin{equation}
  D\kappa^4 = \omega^2 \rho h \left[1 + \frac{\mu\,\omega_r^2}{\omega_r^2 - \omega^2}\right] .
\end{equation}
Bringing the bracket to a common denominator,
\begin{equation}
  1 + \frac{\mu\,\omega_r^2}{\omega_r^2 - \omega^2}
  = \frac{\omega_r^2 - \omega^2 + \mu\,\omega_r^2}{\omega_r^2 - \omega^2}
  = \frac{\omega_r^2(1 + \mu) - \omega^2}{\omega_r^2 - \omega^2},
\end{equation}
so that
\begin{equation}
  D\kappa^4 = \omega^2 \rho h\,\frac{\omega_r^2(1 + \mu) - \omega^2}{\omega_r^2 - \omega^2} .
  \label{eq:effective_medium}
\end{equation}
A real $\kappa$ needs a non-negative right-hand side:
\begin{center}
\begin{tabular}{|l|c|c|l|}
\hline
Range & Numerator & Denominator & $\kappa^4$ \\
\hline
$\omega < \omega_r$ & $+$ & $+$ & $+$: waves propagate \\
$\omega_r < \omega < \omega_r\sqrt{1+\mu}$ & $+$ & $-$ & $-$: no real $\kappa$, gap \\
$\omega > \omega_r\sqrt{1+\mu}$ & $-$ & $-$ & $+$: waves propagate \\
\hline
\end{tabular}
\end{center}
The effective-medium gap is $f_r \dots f_r\sqrt{1+\mu} = 1200 \dots
1342\,\mathrm{Hz}$. The FE result is $1031$--$1195\,\mathrm{Hz}$; the difference
is the topic of the next chapter.

% ---------------------------------------------------------------------
\section{The two-level model with a resonator}

\subsection{The modal amplitude $q$}
Keep one plate mode only, so the deflection is
\begin{equation}
  w(\mathbf{x}) = q\,\varphi(\mathbf{x}),
  \qquad
  \varphi(\mathbf{x}) = \frac{e^{i(\mathbf{k}+\mathbf{G})\cdot\mathbf{x}}}{\sqrt{\rho h A}} .
\end{equation}
$q$ is the \emph{modal amplitude}: how strongly this mode is present. Since
$\varphi$ has units $1/\sqrt{\mathrm{kg}}$, $q$ has units
$\sqrt{\mathrm{kg}}\cdot\mathrm{m}$. If $W$ is the deflection amplitude in
metres, $w = W e^{i(\mathbf{k}+\mathbf{G})\cdot\mathbf{x}}$ and
$q = W\sqrt{\rho h A}$. At the attachment point $w_0 = \varphi\, q$, where
$\varphi$ without an argument means $\varphi(\mathbf{x}_0)$. The unknowns are
collected in the vector $x = (q, u)^T$.

\subsection{Direct route: two force balances}

\paragraph{Row 2: Newton for the resonator.}
The spring stretch is $\varphi q - u$, so $-m\omega^2 u = k_s(\varphi q - u)$.
With $\lambda = \omega^2$:
\begin{equation}
  -k_s\varphi\, q + k_s u = \lambda\, m\, u .
  \label{eq:row_resonator}
\end{equation}

\paragraph{Row 1: the plate mode.}
This is the derivation of Chapter~\ref{ch:two_level} with another force. The
spring pulls on the plate at $\mathbf{x}_0$ with $F = k_s(u - w_0)$:
\begin{equation}
  \mathcal{L} w - \omega^2 \rho h\, w = F\,\delta(\mathbf{x} - \mathbf{x}_0) .
\end{equation}
Insert $w = q\varphi$, multiply by $\varphi^*$ and integrate over the cell. The
$\mathcal{L}$ term gives $\lambda_p q$, the $\rho h$ term $\lambda q$, and the
$\delta$ term $\varphi^* F = \varphi^* k_s(u - \varphi q)$. Hence
$\lambda_p q - \lambda q = k_s\varphi^* u - k_s|\varphi|^2 q$, or
\begin{equation}
  (\lambda_p + k_s|\varphi|^2)\, q - k_s\varphi^*\, u = \lambda\, q .
  \label{eq:row_plate}
\end{equation}

\paragraph{Matrix form.}
Eqs.~\eqref{eq:row_plate} and \eqref{eq:row_resonator} together read
$K x = \lambda M x$ with
\begin{equation}
  K = \begin{pmatrix} \lambda_p + k_s|\varphi|^2 & -k_s\varphi^* \\ -k_s\varphi & k_s \end{pmatrix},
  \qquad
  M = \begin{pmatrix} 1 & 0 \\ 0 & m \end{pmatrix} .
  \label{eq:KM_resonator}
\end{equation}
Each row of $K$ holds the coefficients of $q$ and $u$ on the left of one
equation; each row of $M$ the coefficients on the right. The $1$ in $M$ is the
modal mass of the plate mode, equal to one by the normalization. The entries of
$K$ have different units because $q$ and $u$ do.

\subsection{The same matrices from the energies}

\paragraph{Plate kinetic energy.} Kinetic energy is
$\tfrac12\cdot\text{mass}\cdot\text{velocity}^2$, summed over the plate, with
velocity amplitude $\omega w = \omega q\varphi$:
\begin{equation}
  T_\mathrm{plate} = \tfrac12 \int \rho h\, \omega^2 |q|^2 |\varphi|^2 \,\mathrm{d}A
  = \tfrac12\,\omega^2 |q|^2 \int \rho h\, |\varphi|^2\,\mathrm{d}A
  = \tfrac12\,\omega^2 |q|^2 .
\end{equation}

\paragraph{Resonator kinetic energy.} $T_\mathrm{res} = \tfrac12\, m\, (\omega u)^2 = \tfrac12\,\omega^2 m |u|^2$.

\paragraph{Plate strain energy.} A mode acts like a single oscillator of mass
one. From $\omega^2 = \text{stiffness}/\text{mass}$ its stiffness is
$\lambda_p$, so the energy is $\tfrac12\lambda_p|q|^2$. Formally,
$\tfrac12\int w^*\mathcal{L}w\,\mathrm{d}A = \tfrac12|q|^2\lambda_p\int\rho h|\varphi|^2\,\mathrm{d}A$.

\paragraph{Spring strain energy.} $\tfrac12\, k_s\cdot\text{stretch}^2 = \tfrac12\, k_s|u - \varphi q|^2$.

\paragraph{What multiplying out does.}
Any such energy can be written with a matrix, which is how $K$ and $M$ are
defined:
\begin{align}
  \text{strain energy} &= \tfrac12\, x^H K x = \tfrac12\left(K_{11}|q|^2 + K_{12}\,q^* u + K_{21}\,u^* q + K_{22}|u|^2\right), \\
  \text{kinetic energy} &= \tfrac12\,\omega^2\, x^H M x .
\end{align}
$K$ is read off by sorting the energy into the four products $|q|^2$, $q^*u$,
$u^*q$, $|u|^2$. Multiplying out does that sorting:
\begin{equation}
  |u - \varphi q|^2 = (u^* - \varphi^* q^*)(u - \varphi q)
  = |u|^2 - \varphi\, q\, u^* - \varphi^* q^* u + |\varphi|^2 |q|^2 .
\end{equation}
\begin{center}
\begin{tabular}{|l|l|l|}
\hline
Product & Coefficient in $2\times$ strain energy & Entry \\
\hline
$|q|^2$ & $\lambda_p + k_s|\varphi|^2$ & $K_{11}$ \\
$q^* u$ & $-k_s\varphi^*$ & $K_{12}$ \\
$u^* q$ & $-k_s\varphi$ & $K_{21}$ \\
$|u|^2$ & $k_s$ & $K_{22}$ \\
\hline
\end{tabular}
\end{center}
For $M$, $2\times$ kinetic energy$/\omega^2 = |q|^2 + m|u|^2$ gives
$M_{11} = 1$, $M_{22} = m$. $M$ has no off-diagonal entries because no kinetic
energy term contains both $q$ and $u$: the masses are connected only by the
spring, so the coupling sits in $K$. Energy sorting and force balance are two
routes to the same result; FE assembly uses the energy route.

\paragraph{Link to the code.}
The stretch is $r\,x$ with the row vector $r = (\varphi, -1)$, so the spring
part of $K$ is $k_s\, r^H r$. This is what the \texttt{augment\_*} functions
build: plate $K$ plus this rank-one term, and $m$ on the diagonal of the
resonator DOF. In the code the vector is real (FE basis weights), and the
phases enter through $T(\mathbf{k})$.

\subsection{The standard form $H$}
With the mass-normalized resonator coordinate $v = \sqrt{m}\,u$, $M$ becomes
the identity and $K$ becomes a Hermitian matrix $H$. For a plane wave,
$|\varphi|^2 = 1/(\rho h A)$, and with $\mathbf{x}_0 = 0$, $\varphi$ is real:
\begin{align}
  \text{plate diagonal:} \quad & k_s|\varphi|^2 = \frac{k_s}{\rho h A} = \frac{k_s}{m}\cdot\frac{m}{\rho h A} = \mu\lambda_r, \\
  \text{resonator diagonal:} \quad & \frac{k_s}{m} = \lambda_r, \\
  \text{off-diagonal:} \quad & \frac{k_s\varphi}{\sqrt{m}} = \frac{k_s}{\sqrt{m\,\rho h A}} = \frac{k_s}{m}\sqrt{\frac{m}{\rho h A}} = \sqrt{\mu}\,\lambda_r .
\end{align}
So
\begin{equation}
  \boxed{\;H = \begin{pmatrix} \lambda_p + \mu\lambda_r & -\sqrt{\mu}\,\lambda_r \\ -\sqrt{\mu}\,\lambda_r & \lambda_r \end{pmatrix}\;}
  \label{eq:H_resonator}
\end{equation}
\begin{itemize}
  \item $\lambda_p + \mu\lambda_r$: plate mode, slightly \emph{stiffened} by the
        attached spring;
  \item $\lambda_r$: the resonator alone on a fixed base;
  \item $c = -\sqrt{\mu}\,\lambda_r$: the coupling through the spring.
\end{itemize}
Unlike the rigid-mass matrix, $H$ does not contain the unknown $\lambda$. The
extra DOF makes the problem linear again, which is why the FE code augments
instead of using $m_\mathrm{eff}(\omega)$.

\subsection{Checks}

\paragraph{$\Gamma$, band 1 ($\lambda_p = 0$: the cell translates rigidly).}
$H = \lambda_r \begin{pmatrix} \mu & -\sqrt{\mu} \\ -\sqrt{\mu} & 1 \end{pmatrix}$
has trace $\lambda_r(1+\mu)$ and determinant $\lambda_r^2(\mu - \mu) = 0$. The
eigenvalues are $0$ (plate and resonator translate together) and
$\lambda_r(1+\mu)$, the upper edge of the effective-medium gap.

\paragraph{Short waves ($\lambda_p \gg \lambda_r$).}
Expanding as in Section~\ref{sec:far}, the lower branch tends to $\lambda_r$:
the plate is too stiff to move, and the resonator vibrates on a fixed base.

\paragraph{Back to the effective medium.}
$\det(H - \lambda I) = 0$ reads
\begin{equation}
  (\lambda_p + \mu\lambda_r - \lambda)(\lambda_r - \lambda) - \mu\lambda_r^2 = 0 .
\end{equation}
Expanding the product,
\begin{equation}
  \lambda_p(\lambda_r - \lambda) + \mu\lambda_r(\lambda_r - \lambda) - \lambda(\lambda_r - \lambda) - \mu\lambda_r^2 = 0 .
\end{equation}
The $\mu\lambda_r^2$ terms cancel:
\begin{equation}
  \lambda_p(\lambda_r - \lambda) - \mu\lambda_r\lambda - \lambda(\lambda_r - \lambda) = 0 .
\end{equation}
Solving for $\lambda_p$:
\begin{equation}
  \lambda_p = \lambda\,\frac{\lambda_r(1+\mu) - \lambda}{\lambda_r - \lambda} .
\end{equation}
With $\lambda_p = D\kappa^4/\rho h$ this is exactly
Eq.~\eqref{eq:effective_medium}. The two-level model with one plate mode
\emph{is} the effective medium.

\subsection{The avoided crossing}
The diagonal entries are equal where $\lambda_p + \mu\lambda_r = \lambda_r$,
i.e.\ $\lambda_p = \lambda_r(1-\mu)$. There the splitting is $\pm|c|$:
\begin{equation}
  \lambda = \lambda_r \pm \sqrt{\mu}\,\lambda_r = \lambda_r(1 \pm \sqrt{\mu}) .
\end{equation}
For $\mu = 0.25$: $f = f_r\sqrt{0.5} \approx 849\,\mathrm{Hz}$ and
$f = f_r\sqrt{1.5} \approx 1470\,\mathrm{Hz}$.

\paragraph{Splitting versus gap width.}
$\pm\sqrt{\mu}$ is the spacing of the two bands \emph{at that one
$\mathbf{k}$}. In the unfolded model the complete gap runs from the
\emph{maximum} of the lower band (tending to $f_r$) to the \emph{minimum} of the
upper band ($f_r\sqrt{1+\mu}$ at $\Gamma$). Splitting and gap width are
different things.

\paragraph{Why $\sqrt{\mu}$ here and $\mu$ for the rigid mass?}
In Chapter~\ref{ch:two_level} two \emph{plate} modes were coupled indirectly
through the inertia of the mass; both ends of the chain carried a factor
$m|\varphi|^2$-type term, so the coupling went as $\mu$. Here the resonator is
a degree of freedom of its own, attached directly by the spring. Each
normalized coordinate carries one square root of its mass, which gives
$\sqrt{\mu}$.

% ---------------------------------------------------------------------
\section{The flat band and the plate share}

\paragraph{Where the flat band comes from.}
The augmented problem has one more DOF per cell than the bare plate, so there
is one extra eigenvalue at every $\mathbf{k}$. Away from plate bands it is the
resonator vibrating while the plate moves little. The resonator is a
\emph{local} DOF, and $\mathbf{k}$ enters only through the phase factors of
the plate; if the plate moves little, the eigenvalue barely depends on
$\mathbf{k}$.

\paragraph{The slope of a band.}
$H$ in Eq.~\eqref{eq:H_resonator} depends on $\mathbf{k}$ only through
$\lambda_p$, in one entry. First-order perturbation theory then gives, for a
normalized eigenvector $(q, v)$,
\begin{equation}
  \frac{\mathrm{d}\lambda}{\mathrm{d}\lambda_p} = |q|^2 = \text{plate share of the mode},
\end{equation}
where \emph{plate share} means the fraction of the kinetic energy of the mode
that is in the plate. Therefore
\begin{equation}
  \text{slope of a band} = \text{plate share} \times \text{slope of the plate curve} .
\end{equation}
A band is flat where the plate share is small, or where the plate curve itself
is flat, and it bends where the plate share grows. This happens near the plate
curve, not simply at some distance from $\omega_r$.

\paragraph{Computing the plate share.}
Row 2 of Eq.~\eqref{eq:H_resonator} gives
$-\sqrt{\mu}\,\lambda_r\, q + \lambda_r v = \lambda v$, hence
\begin{equation}
  \frac{q}{v} = \frac{\lambda_r - \lambda}{\sqrt{\mu}\,\lambda_r},
  \qquad
  \text{plate share} = \frac{(q/v)^2}{1 + (q/v)^2} .
\end{equation}
At $X$ and $M$ the coupled combination has $|\varphi|^2 = 2/(\rho h A)$ and
$4/(\rho h A)$ (Chapter~\ref{ch:bragg_2d}), so in all formulas $\mu$ is
replaced by $2\mu$ and $4\mu$, and $\lambda_p$ by the bare value
($0.983\,\lambda_r$ at $X$, $3.93\,\lambda_r$ at $M$).

% ---------------------------------------------------------------------
\section{Reading the dispersion diagram}

\begin{figure}[htbp]
  \centering
  % \includegraphics[width=0.85\textwidth]{dispersion_vamm_1200.png}
  \caption{Dispersion diagram of the $0.1\,\mathrm{m}$ unit cell
    ($h = 5\,\mathrm{mm}$, steel) with one resonator per cell,
    $f_r = 1200\,\mathrm{Hz}$, $\mu = 0.25$, along
    $\Gamma$--$X$--$M$--$\Gamma$. Blue dots: FE Bloch solution. Solid grey:
    plane-wave expansion. Dashed grey: bare plate. Dotted red: resonator
    frequency $f_r$. The complete gap runs from the maximum of band 1 at $M$
    ($1031\,\mathrm{Hz}$) to the minimum of band 2 at $X$
    ($1195\,\mathrm{Hz}$).}
  \label{fig:dispersion_vamm_1200}
\end{figure}

The blue dots in Fig.~\ref{fig:dispersion_vamm_1200} are the actual modes. They
come from three ingredients:
\begin{enumerate}
  \item the grey dashed rising curve: the plate wave going out from $\Gamma$;
  \item the grey dashed falling curves: the folded plate waves coming down
        towards $X$ and $M$;
  \item the red dotted line: the resonator alone, at $1200\,\mathrm{Hz}$.
\end{enumerate}
Wherever these meet, the blue bands repel. Between meetings, each blue band
follows whichever ingredient is closest.

The cleanest place to see this is the right part of the plot, $M$--$\Gamma$.
The rising plate curve crosses $1200\,\mathrm{Hz}$ about $70\,\%$ of the way
from $\Gamma$ to $M$ (at $\kappa = \pi/a$); the diagonal entries of $H$ are equal
slightly earlier, at about two thirds of the way. Reading from $\Gamma$ towards
$M$:
\begin{itemize}
  \item \textbf{Band 1} starts on the grey curve, bends away below the crossing
        and flattens out.
  \item \textbf{Band 2} starts flat at about $1270\,\mathrm{Hz}$, stays flat,
        then bends upwards at the crossing and joins the grey curve.
\end{itemize}
The two bands swap roles at the crossing: band 1 starts plate-like and ends
resonator-like, band 2 does the opposite. This is the avoided crossing of the
previous section, and its model values $849$ and $1470\,\mathrm{Hz}$ fit the
plot at that position.

\paragraph{Why the bands are flat where they are.}
\begin{itemize}
  \item Flat at $1270\,\mathrm{Hz}$ near $\Gamma$: the plate share is low
        ($20\,\%$ in the model) and the plate curve itself is nearly flat
        there. Both reasons work together.
  \item Band 1 near $\Gamma$: the share is high ($80\,\%$), so the band follows
        the grey curve.
  \item The bending part of band 2 on $M$--$\Gamma$: the share grows towards
        $100\,\%$, and the band picks up the slope of the grey curve.
\end{itemize}

\paragraph{Plate share at the landmarks.}
The table extends the $2\times 2$ model to the coupled combinations at $X$ and
$M$; ``crossing'' is the avoided crossing on $M$--$\Gamma$. The shares are \emph{model estimates, not FE results}. The FE value
$2710\,\mathrm{Hz}$ was read off the plot.
\begin{center}
\begin{tabular}{|l|l|l|l|p{0.26\textwidth}|}
\hline
Point & Band & FE [Hz] & Model [Hz] & Plate share \\
\hline
$\Gamma$ & 1 & $0$ & $0$ & $80\,\%$ (resonators ride along) \\
$\Gamma$ & 2 (flat) & $\approx 1270$ & $1342$ & $20\,\%$ \\
crossing & 1 / 2 & -- & $849$ / $1470$ & $50\,\%$ / $50\,\%$ \\
$X$ & 1 & $836$ & $846$ & $34\,\%$ \\
$X$ & 2 & $1195$ & $1194$ & $100\,\%$ (node at the resonator) \\
$X$ & 3 & $1652$ & $1694$ & $66\,\%$ \\
$M$ & 1 & $1031$ & $1046$ & $5\,\%$ \\
$M$ & top of group & $\approx 2710$ & $2728$ & $95\,\%$ \\
\hline
\end{tabular}
\end{center}

\paragraph{Two points that are easy to get wrong.}
\begin{enumerate}
  \item Band 1 never reaches $f_r$ in this diagram, because the zone ends first.
        It tops out at $1031\,\mathrm{Hz}$ at $M$. The limit $f \to f_r$ exists
        only in the unfolded model; in the diagram it corresponds to band 1
        near $M$, with a plate share of $5\,\%$.
  \item The ``upper band at $\Gamma$'' of the two-level model \emph{is} the flat
        band of the plot. It sits at about $1270\,\mathrm{Hz}$, above the red
        line. Nothing in the diagram sits on $1200\,\mathrm{Hz}$.
\end{enumerate}

\paragraph{Outlook.}
The gap runs from the top of band 1 at $M$ to band 2 at $X$. That upper edge is
the $100\,\%$ plate mode, which does not feel the resonator at all. The next
chapter takes this up.

% ---------------------------------------------------------------------
\section{Check questions with answers}

\paragraph{Question 1.} \emph{Just above $f_r$, $m_\mathrm{eff}$ is large and
negative. What does the plate feel there, and why does that rule out a
propagating bending wave?}

\emph{Answer.} The resonator oscillates opposite to the plate and pulls it back,
so the usual propagating wave is stopped. Two refinements:
\begin{itemize}
  \item Nothing is dissipated, because the model is undamped. The wave becomes
        \emph{evanescent}: $\kappa^4 < 0$ makes $\kappa$ complex, so the plate
        still moves, but the motion decays from cell to cell. The energy is
        reflected, not absorbed.
  \item Negative $m_\mathrm{eff}$ alone is not enough. $m_\mathrm{eff}$ is
        negative at \emph{every} frequency above $f_r$, but the gap ends at
        $f_r\sqrt{1+\mu}$. What counts is the \emph{total} mass per area,
        $\rho h + m_\mathrm{eff}/A$. It is negative only while the negative
        mass of the resonator outweighs the mass of the plate, and it crosses
        zero exactly at $\omega^2 = \omega_r^2(1+\mu)$.
\end{itemize}

\paragraph{Question 2.} \emph{Why is the resonator band flat, and where would you
expect it not to be flat?}

\emph{Answer.} The plate moves little there, so the frequency barely depends on
$\mathbf{k}$. What decides flatness is not the distance from $\omega_r$ itself,
but how far the \emph{plate band at that $\mathbf{k}$} is from the resonator
level: slope of a band = plate share $\times$ slope of the plate curve. The
band bends at the avoided crossing, and it moves away from $\omega_r$
\emph{because} it hybridizes. The band near $\Gamma$ shows the difference: it
sits at a frequency different from $f_r$ ($1342\,\mathrm{Hz}$ in the model,
about $1270\,\mathrm{Hz}$ in FE) and is still nearly flat, because the plate
band is far below it there.

\paragraph{Question 3.} \emph{Along the lower band, the eigenvector goes from
mostly plate to mostly resonator. What does this mean for MAC tracking, and
which branch does a shaker on the finite plate see?}

\emph{Answer.} The 50/50 mix sits at the crossing, not near $f_r$; the mode
becomes resonator-dominated towards the end of band 1.

\emph{MAC tracking:} band 1 goes from $80\,\%$ plate at $\Gamma$ to $5\,\%$ at
$M$. That is a large change of character, but it happens gradually over many
$\mathbf{k}$-points, so the MAC between neighbouring points stays high and
tracking should follow the band without swapping labels. The MAC must use the
\emph{full vector including the resonator DOF}, with the augmented mass matrix;
on plate DOFs alone the comparison is somewhat weaker (lowest MAC of band 1:
$0.969$ instead of $0.990$), and leaving out the mass matrix is far worse ($0.619$).

\emph{Shaker:} the question has no clean one-branch answer. The shaker drives the
plate, so a mode is excited in proportion to its plate share. It sees band 1 well
below about $850\,\mathrm{Hz}$, band 1 near $M$ hardly at all, and the flat part
of band 2 only weakly (reasoning from the model, not measured). Near the gap
edges the upper branch keeps more plate motion than the lower one: at the upper
edge at $\Gamma$ the model gives a plate share of $\mu/(1+\mu) = 20\,\%$ and
$u/w = -1/\mu = -4$.

% =====================================================================

\chapter{The local-resonance gap}
\label{ch:local_resonance}

This chapter replaces the rigid mass of Chapters~\ref{ch:two_level} and
\ref{ch:bragg_2d} by a real mass--spring resonator. The structure stays the
same: one attachment point per cell, coupling only through
$w(\mathbf{x}_0)$, and two-level models.

\begin{center}
\begin{tabular}{|l|p{0.45\textwidth}|l|}
\hline
Symbol & Meaning & Value / units \\
\hline
$m$ & resonator mass, one per cell & $98\,\mathrm{g}$ \\
$k_s$ & spring stiffness (subscript $s$ to avoid confusion with the wavevector $\mathbf{k}$) & $\approx 5.6\cdot 10^6\,\mathrm{N/m}$ \\
$\lambda_r = \omega_r^2 = k_s/m$ & squared angular tuning frequency & $f_r = 1200\,\mathrm{Hz}$ \\
$u$ & displacement of the resonator mass & m \\
$w_0 = w(\mathbf{x}_0)$ & plate deflection at the attachment point & m \\
$\mu = m/(\rho h A)$ & mass ratio & $0.25$ \\
$\lambda = \omega^2$ & \textbf{unknown}: $\omega^2$ of a mode of plate and resonator (one dot of the diagram) & $\mathrm{rad^2/s^2}$ \\
$\lambda_p(\mathbf{k})$ & $\omega^2$ of one bare plate mode at this $\mathbf{k}$ (the $\lambda_n^0$ of Chapter~\ref{ch:two_level}) & $\mathrm{rad^2/s^2}$ \\
$q$ & modal amplitude of that plate mode & $\sqrt{\mathrm{kg}}\cdot\mathrm{m}$ \\
\hline
\end{tabular}
\end{center}

% ---------------------------------------------------------------------
\section{The resonator as a frequency-dependent mass}

\paragraph{Equation of motion.}
The spring is stretched by $w_0 - u$, so $m\ddot{u} = k_s(w_0 - u)$. For
harmonic motion, $\ddot{u} = -\omega^2 u$:
\begin{align}
  -m\omega^2 u &= k_s w_0 - k_s u \\
  u\,(k_s - m\omega^2) &= k_s w_0 \\
  u &= \frac{k_s w_0}{k_s - m\omega^2} = w_0\,\frac{\omega_r^2}{\omega_r^2 - \omega^2},
\end{align}
where the last step divides numerator and denominator by $m$.

\paragraph{Force on the plate.}
The spring pulls the plate towards the resonator, $F = k_s(u - w_0)$. With
\begin{equation}
  u - w_0 = w_0\left[\frac{\omega_r^2}{\omega_r^2 - \omega^2} - 1\right]
          = w_0\,\frac{\omega^2}{\omega_r^2 - \omega^2}
\end{equation}
and $k_s = m\omega_r^2$ we get
\begin{equation}
  F = \omega^2 w_0 \cdot \frac{m\,\omega_r^2}{\omega_r^2 - \omega^2} .
\end{equation}
A rigid mass pushed on the plate with $F = m\omega^2 w_0$. The resonator gives
the same expression with $m$ replaced by the \emph{effective mass}
\begin{equation}
  \boxed{\;m_\mathrm{eff}(\omega) = \frac{m\,\omega_r^2}{\omega_r^2 - \omega^2}\;}
\end{equation}
Everything from the rigid-mass chapters carries over by replacing $m$ with
$m_\mathrm{eff}(\omega)$ and $\mu$ with
$\mu_\mathrm{eff}(\omega) = \mu\,\omega_r^2/(\omega_r^2 - \omega^2)$.

\begin{center}
\begin{tabular}{|l|l|p{0.22\textwidth}|p{0.3\textwidth}|}
\hline
Frequency & $m_\mathrm{eff}$ & Resonator motion & What the plate feels \\
\hline
$\omega \ll \omega_r$ & $\approx m$ & $u \approx w_0$, with the plate & a rigid mass \\
just below $\omega_r$ & $\to +\infty$ & $u \gg w_0$, in phase & a huge mass: the plate is pinned \\
just above $\omega_r$ & $\to -\infty$ & $u \gg w_0$, antiphase & the spring pulls the plate back against its motion \\
$\omega \gg \omega_r$ & $\approx -m\omega_r^2/\omega^2$ & $u \approx 0$, mass stands still & $F \to -k_s w_0$: a spring to a fixed point \\
\hline
\end{tabular}
\end{center}
A negative $m_\mathrm{eff}$ is not strange. It means that the force of the
resonator acts as a \emph{restoring} force (stiffness) and not as an inertial
one. The last row shows this most clearly.

% ---------------------------------------------------------------------
\section{Effective medium: smearing the resonators over the plate}

For waves much longer than $a$ only $\mathbf{G} = 0$ matters, and $m_\mathrm{eff}$
can be spread over the cell area: $\rho h \to \rho h + m_\mathrm{eff}/A$. The
Kirchhoff relation $D\kappa^4 = \rho h\,\omega^2$ becomes
\begin{equation}
  D\kappa^4 = \omega^2 \rho h \left[1 + \frac{\mu\,\omega_r^2}{\omega_r^2 - \omega^2}\right] .
\end{equation}
Bringing the bracket to a common denominator,
\begin{equation}
  1 + \frac{\mu\,\omega_r^2}{\omega_r^2 - \omega^2}
  = \frac{\omega_r^2 - \omega^2 + \mu\,\omega_r^2}{\omega_r^2 - \omega^2}
  = \frac{\omega_r^2(1 + \mu) - \omega^2}{\omega_r^2 - \omega^2},
\end{equation}
so that
\begin{equation}
  D\kappa^4 = \omega^2 \rho h\,\frac{\omega_r^2(1 + \mu) - \omega^2}{\omega_r^2 - \omega^2} .
  \label{eq:effective_medium}
\end{equation}
A real $\kappa$ needs a non-negative right-hand side:
\begin{center}
\begin{tabular}{|l|c|c|l|}
\hline
Range & Numerator & Denominator & $\kappa^4$ \\
\hline
$\omega < \omega_r$ & $+$ & $+$ & $+$: waves propagate \\
$\omega_r < \omega < \omega_r\sqrt{1+\mu}$ & $+$ & $-$ & $-$: no real $\kappa$, gap \\
$\omega > \omega_r\sqrt{1+\mu}$ & $-$ & $-$ & $+$: waves propagate \\
\hline
\end{tabular}
\end{center}
The effective-medium gap is $f_r \dots f_r\sqrt{1+\mu} = 1200 \dots
1342\,\mathrm{Hz}$. The FE result is $1031$--$1195\,\mathrm{Hz}$; the difference
is the topic of the next chapter.

% ---------------------------------------------------------------------
\section{The two-level model with a resonator}

\subsection{The modal amplitude $q$}
Keep one plate mode only, so the deflection is
\begin{equation}
  w(\mathbf{x}) = q\,\varphi(\mathbf{x}),
  \qquad
  \varphi(\mathbf{x}) = \frac{e^{i(\mathbf{k}+\mathbf{G})\cdot\mathbf{x}}}{\sqrt{\rho h A}} .
\end{equation}
$q$ is the \emph{modal amplitude}: how strongly this mode is present. Since
$\varphi$ has units $1/\sqrt{\mathrm{kg}}$, $q$ has units
$\sqrt{\mathrm{kg}}\cdot\mathrm{m}$. If $W$ is the deflection amplitude in
metres, $w = W e^{i(\mathbf{k}+\mathbf{G})\cdot\mathbf{x}}$ and
$q = W\sqrt{\rho h A}$. At the attachment point $w_0 = \varphi\, q$, where
$\varphi$ without an argument means $\varphi(\mathbf{x}_0)$. The unknowns are
collected in the vector $x = (q, u)^T$.

\subsection{Direct route: two force balances}

\paragraph{Row 2: Newton for the resonator.}
The spring stretch is $\varphi q - u$, so $-m\omega^2 u = k_s(\varphi q - u)$.
With $\lambda = \omega^2$:
\begin{equation}
  -k_s\varphi\, q + k_s u = \lambda\, m\, u .
  \label{eq:row_resonator}
\end{equation}

\paragraph{Row 1: the plate mode.}
This is the derivation of Chapter~\ref{ch:two_level} with another force. The
spring pulls on the plate at $\mathbf{x}_0$ with $F = k_s(u - w_0)$:
\begin{equation}
  \mathcal{L} w - \omega^2 \rho h\, w = F\,\delta(\mathbf{x} - \mathbf{x}_0) .
\end{equation}
Insert $w = q\varphi$, multiply by $\varphi^*$ and integrate over the cell. The
$\mathcal{L}$ term gives $\lambda_p q$, the $\rho h$ term $\lambda q$, and the
$\delta$ term $\varphi^* F = \varphi^* k_s(u - \varphi q)$. Hence
$\lambda_p q - \lambda q = k_s\varphi^* u - k_s|\varphi|^2 q$, or
\begin{equation}
  (\lambda_p + k_s|\varphi|^2)\, q - k_s\varphi^*\, u = \lambda\, q .
  \label{eq:row_plate}
\end{equation}

\paragraph{Matrix form.}
Eqs.~\eqref{eq:row_plate} and \eqref{eq:row_resonator} together read
$K x = \lambda M x$ with
\begin{equation}
  K = \begin{pmatrix} \lambda_p + k_s|\varphi|^2 & -k_s\varphi^* \\ -k_s\varphi & k_s \end{pmatrix},
  \qquad
  M = \begin{pmatrix} 1 & 0 \\ 0 & m \end{pmatrix} .
  \label{eq:KM_resonator}
\end{equation}
Each row of $K$ holds the coefficients of $q$ and $u$ on the left of one
equation; each row of $M$ the coefficients on the right. The $1$ in $M$ is the
modal mass of the plate mode, equal to one by the normalization. The entries of
$K$ have different units because $q$ and $u$ do.

\subsection{The same matrices from the energies}

\paragraph{Plate kinetic energy.} Kinetic energy is
$\tfrac12\cdot\text{mass}\cdot\text{velocity}^2$, summed over the plate, with
velocity amplitude $\omega w = \omega q\varphi$:
\begin{equation}
  T_\mathrm{plate} = \tfrac12 \int \rho h\, \omega^2 |q|^2 |\varphi|^2 \,\mathrm{d}A
  = \tfrac12\,\omega^2 |q|^2 \int \rho h\, |\varphi|^2\,\mathrm{d}A
  = \tfrac12\,\omega^2 |q|^2 .
\end{equation}

\paragraph{Resonator kinetic energy.} $T_\mathrm{res} = \tfrac12\, m\, (\omega u)^2 = \tfrac12\,\omega^2 m |u|^2$.

\paragraph{Plate strain energy.} A mode acts like a single oscillator of mass
one. From $\omega^2 = \text{stiffness}/\text{mass}$ its stiffness is
$\lambda_p$, so the energy is $\tfrac12\lambda_p|q|^2$. Formally,
$\tfrac12\int w^*\mathcal{L}w\,\mathrm{d}A = \tfrac12|q|^2\lambda_p\int\rho h|\varphi|^2\,\mathrm{d}A$.

\paragraph{Spring strain energy.} $\tfrac12\, k_s\cdot\text{stretch}^2 = \tfrac12\, k_s|u - \varphi q|^2$.

\paragraph{What multiplying out does.}
Any such energy can be written with a matrix, which is how $K$ and $M$ are
defined:
\begin{align}
  \text{strain energy} &= \tfrac12\, x^H K x = \tfrac12\left(K_{11}|q|^2 + K_{12}\,q^* u + K_{21}\,u^* q + K_{22}|u|^2\right), \\
  \text{kinetic energy} &= \tfrac12\,\omega^2\, x^H M x .
\end{align}
$K$ is read off by sorting the energy into the four products $|q|^2$, $q^*u$,
$u^*q$, $|u|^2$. Multiplying out does that sorting:
\begin{equation}
  |u - \varphi q|^2 = (u^* - \varphi^* q^*)(u - \varphi q)
  = |u|^2 - \varphi\, q\, u^* - \varphi^* q^* u + |\varphi|^2 |q|^2 .
\end{equation}
\begin{center}
\begin{tabular}{|l|l|l|}
\hline
Product & Coefficient in $2\times$ strain energy & Entry \\
\hline
$|q|^2$ & $\lambda_p + k_s|\varphi|^2$ & $K_{11}$ \\
$q^* u$ & $-k_s\varphi^*$ & $K_{12}$ \\
$u^* q$ & $-k_s\varphi$ & $K_{21}$ \\
$|u|^2$ & $k_s$ & $K_{22}$ \\
\hline
\end{tabular}
\end{center}
For $M$, $2\times$ kinetic energy$/\omega^2 = |q|^2 + m|u|^2$ gives
$M_{11} = 1$, $M_{22} = m$. $M$ has no off-diagonal entries because no kinetic
energy term contains both $q$ and $u$: the masses are connected only by the
spring, so the coupling sits in $K$. Energy sorting and force balance are two
routes to the same result; FE assembly uses the energy route.

\paragraph{Link to the code.}
The stretch is $r\,x$ with the row vector $r = (\varphi, -1)$, so the spring
part of $K$ is $k_s\, r^H r$. This is what the \texttt{augment\_*} functions
build: plate $K$ plus this rank-one term, and $m$ on the diagonal of the
resonator DOF. In the code the vector is real (FE basis weights), and the
phases enter through $T(\mathbf{k})$.

\subsection{The standard form $H$}
With the mass-normalized resonator coordinate $v = \sqrt{m}\,u$, $M$ becomes
the identity and $K$ becomes a Hermitian matrix $H$. For a plane wave,
$|\varphi|^2 = 1/(\rho h A)$, and with $\mathbf{x}_0 = 0$, $\varphi$ is real:
\begin{align}
  \text{plate diagonal:} \quad & k_s|\varphi|^2 = \frac{k_s}{\rho h A} = \frac{k_s}{m}\cdot\frac{m}{\rho h A} = \mu\lambda_r, \\
  \text{resonator diagonal:} \quad & \frac{k_s}{m} = \lambda_r, \\
  \text{off-diagonal:} \quad & \frac{k_s\varphi}{\sqrt{m}} = \frac{k_s}{\sqrt{m\,\rho h A}} = \frac{k_s}{m}\sqrt{\frac{m}{\rho h A}} = \sqrt{\mu}\,\lambda_r .
\end{align}
So
\begin{equation}
  \boxed{\;H = \begin{pmatrix} \lambda_p + \mu\lambda_r & -\sqrt{\mu}\,\lambda_r \\ -\sqrt{\mu}\,\lambda_r & \lambda_r \end{pmatrix}\;}
  \label{eq:H_resonator}
\end{equation}
\begin{itemize}
  \item $\lambda_p + \mu\lambda_r$: plate mode, slightly \emph{stiffened} by the
        attached spring;
  \item $\lambda_r$: the resonator alone on a fixed base;
  \item $c = -\sqrt{\mu}\,\lambda_r$: the coupling through the spring.
\end{itemize}
Unlike the rigid-mass matrix, $H$ does not contain the unknown $\lambda$. The
extra DOF makes the problem linear again, which is why the FE code augments
instead of using $m_\mathrm{eff}(\omega)$.

\subsection{Checks}

\paragraph{$\Gamma$, band 1 ($\lambda_p = 0$: the cell translates rigidly).}
$H = \lambda_r \begin{pmatrix} \mu & -\sqrt{\mu} \\ -\sqrt{\mu} & 1 \end{pmatrix}$
has trace $\lambda_r(1+\mu)$ and determinant $\lambda_r^2(\mu - \mu) = 0$. The
eigenvalues are $0$ (plate and resonator translate together) and
$\lambda_r(1+\mu)$, the upper edge of the effective-medium gap.

\paragraph{Short waves ($\lambda_p \gg \lambda_r$).}
Expanding as in Section~\ref{sec:far}, the lower branch tends to $\lambda_r$:
the plate is too stiff to move, and the resonator vibrates on a fixed base.

\paragraph{Back to the effective medium.}
$\det(H - \lambda I) = 0$ reads
\begin{equation}
  (\lambda_p + \mu\lambda_r - \lambda)(\lambda_r - \lambda) - \mu\lambda_r^2 = 0 .
\end{equation}
Expanding the product,
\begin{equation}
  \lambda_p(\lambda_r - \lambda) + \mu\lambda_r(\lambda_r - \lambda) - \lambda(\lambda_r - \lambda) - \mu\lambda_r^2 = 0 .
\end{equation}
The $\mu\lambda_r^2$ terms cancel:
\begin{equation}
  \lambda_p(\lambda_r - \lambda) - \mu\lambda_r\lambda - \lambda(\lambda_r - \lambda) = 0 .
\end{equation}
Solving for $\lambda_p$:
\begin{equation}
  \lambda_p = \lambda\,\frac{\lambda_r(1+\mu) - \lambda}{\lambda_r - \lambda} .
\end{equation}
With $\lambda_p = D\kappa^4/\rho h$ this is exactly
Eq.~\eqref{eq:effective_medium}. The two-level model with one plate mode
\emph{is} the effective medium.

\subsection{The avoided crossing}
The diagonal entries are equal where $\lambda_p + \mu\lambda_r = \lambda_r$,
i.e.\ $\lambda_p = \lambda_r(1-\mu)$. There the splitting is $\pm|c|$:
\begin{equation}
  \lambda = \lambda_r \pm \sqrt{\mu}\,\lambda_r = \lambda_r(1 \pm \sqrt{\mu}) .
\end{equation}
For $\mu = 0.25$: $f = f_r\sqrt{0.5} \approx 849\,\mathrm{Hz}$ and
$f = f_r\sqrt{1.5} \approx 1470\,\mathrm{Hz}$.

\paragraph{Splitting versus gap width.}
$\pm\sqrt{\mu}$ is the spacing of the two bands \emph{at that one
$\mathbf{k}$}. In the unfolded model the complete gap runs from the
\emph{maximum} of the lower band (tending to $f_r$) to the \emph{minimum} of the
upper band ($f_r\sqrt{1+\mu}$ at $\Gamma$). Splitting and gap width are
different things.

\paragraph{Why $\sqrt{\mu}$ here and $\mu$ for the rigid mass?}
In Chapter~\ref{ch:two_level} two \emph{plate} modes were coupled indirectly
through the inertia of the mass; both ends of the chain carried a factor
$m|\varphi|^2$-type term, so the coupling went as $\mu$. Here the resonator is
a degree of freedom of its own, attached directly by the spring. Each
normalized coordinate carries one square root of its mass, which gives
$\sqrt{\mu}$.

% ---------------------------------------------------------------------
\section{The flat band and the plate share}

\paragraph{Where the flat band comes from.}
The augmented problem has one more DOF per cell than the bare plate, so there
is one extra eigenvalue at every $\mathbf{k}$. Away from plate bands it is the
resonator vibrating while the plate moves little. The resonator is a
\emph{local} DOF, and $\mathbf{k}$ enters only through the phase factors of
the plate; if the plate moves little, the eigenvalue barely depends on
$\mathbf{k}$.

\paragraph{The slope of a band.}
$H$ in Eq.~\eqref{eq:H_resonator} depends on $\mathbf{k}$ only through
$\lambda_p$, in one entry. First-order perturbation theory then gives, for a
normalized eigenvector $(q, v)$,
\begin{equation}
  \frac{\mathrm{d}\lambda}{\mathrm{d}\lambda_p} = |q|^2 = \text{plate share of the mode},
\end{equation}
where \emph{plate share} means the fraction of the kinetic energy of the mode
that is in the plate. Therefore
\begin{equation}
  \text{slope of a band} = \text{plate share} \times \text{slope of the plate curve} .
\end{equation}
A band is flat where the plate share is small, or where the plate curve itself
is flat, and it bends where the plate share grows. This happens near the plate
curve, not simply at some distance from $\omega_r$.

\paragraph{Computing the plate share.}
Row 2 of Eq.~\eqref{eq:H_resonator} gives
$-\sqrt{\mu}\,\lambda_r\, q + \lambda_r v = \lambda v$, hence
\begin{equation}
  \frac{q}{v} = \frac{\lambda_r - \lambda}{\sqrt{\mu}\,\lambda_r},
  \qquad
  \text{plate share} = \frac{(q/v)^2}{1 + (q/v)^2} .
\end{equation}
At $X$ and $M$ the coupled combination has $|\varphi|^2 = 2/(\rho h A)$ and
$4/(\rho h A)$ (Chapter~\ref{ch:bragg_2d}), so in all formulas $\mu$ is
replaced by $2\mu$ and $4\mu$, and $\lambda_p$ by the bare value
($0.983\,\lambda_r$ at $X$, $3.93\,\lambda_r$ at $M$).

% ---------------------------------------------------------------------
\section{Reading the dispersion diagram}

\begin{figure}[htbp]
  \centering
  % \includegraphics[width=0.85\textwidth]{dispersion_vamm_1200.png}
  \caption{Dispersion diagram of the $0.1\,\mathrm{m}$ unit cell
    ($h = 5\,\mathrm{mm}$, steel) with one resonator per cell,
    $f_r = 1200\,\mathrm{Hz}$, $\mu = 0.25$, along
    $\Gamma$--$X$--$M$--$\Gamma$. Blue dots: FE Bloch solution. Solid grey:
    plane-wave expansion. Dashed grey: bare plate. Dotted red: resonator
    frequency $f_r$. The complete gap runs from the maximum of band 1 at $M$
    ($1031\,\mathrm{Hz}$) to the minimum of band 2 at $X$
    ($1195\,\mathrm{Hz}$).}
  \label{fig:dispersion_vamm_1200}
\end{figure}

The blue dots in Fig.~\ref{fig:dispersion_vamm_1200} are the actual modes. They
come from three ingredients:
\begin{enumerate}
  \item the grey dashed rising curve: the plate wave going out from $\Gamma$;
  \item the grey dashed falling curves: the folded plate waves coming down
        towards $X$ and $M$;
  \item the red dotted line: the resonator alone, at $1200\,\mathrm{Hz}$.
\end{enumerate}
Wherever these meet, the blue bands repel. Between meetings, each blue band
follows whichever ingredient is closest.

The cleanest place to see this is the right part of the plot, $M$--$\Gamma$.
The rising plate curve crosses $1200\,\mathrm{Hz}$ about $70\,\%$ of the way
from $\Gamma$ to $M$ (at $\kappa = \pi/a$); the diagonal entries of $H$ are equal
slightly earlier, at about two thirds of the way. Reading from $\Gamma$ towards
$M$:
\begin{itemize}
  \item \textbf{Band 1} starts on the grey curve, bends away below the crossing
        and flattens out.
  \item \textbf{Band 2} starts flat at about $1270\,\mathrm{Hz}$, stays flat,
        then bends upwards at the crossing and joins the grey curve.
\end{itemize}
The two bands swap roles at the crossing: band 1 starts plate-like and ends
resonator-like, band 2 does the opposite. This is the avoided crossing of the
previous section, and its model values $849$ and $1470\,\mathrm{Hz}$ fit the
plot at that position.

\paragraph{Why the bands are flat where they are.}
\begin{itemize}
  \item Flat at $1270\,\mathrm{Hz}$ near $\Gamma$: the plate share is low
        ($20\,\%$ in the model) and the plate curve itself is nearly flat
        there. Both reasons work together.
  \item Band 1 near $\Gamma$: the share is high ($80\,\%$), so the band follows
        the grey curve.
  \item The bending part of band 2 on $M$--$\Gamma$: the share grows towards
        $100\,\%$, and the band picks up the slope of the grey curve.
\end{itemize}

\paragraph{Plate share at the landmarks.}
The table extends the $2\times 2$ model to the coupled combinations at $X$ and
$M$; ``crossing'' is the avoided crossing on $M$--$\Gamma$. The shares are \emph{model estimates, not FE results}. The FE value
$2710\,\mathrm{Hz}$ was read off the plot.
\begin{center}
\begin{tabular}{|l|l|l|l|p{0.26\textwidth}|}
\hline
Point & Band & FE [Hz] & Model [Hz] & Plate share \\
\hline
$\Gamma$ & 1 & $0$ & $0$ & $80\,\%$ (resonators ride along) \\
$\Gamma$ & 2 (flat) & $\approx 1270$ & $1342$ & $20\,\%$ \\
crossing & 1 / 2 & -- & $849$ / $1470$ & $50\,\%$ / $50\,\%$ \\
$X$ & 1 & $836$ & $846$ & $34\,\%$ \\
$X$ & 2 & $1195$ & $1194$ & $100\,\%$ (node at the resonator) \\
$X$ & 3 & $1652$ & $1694$ & $66\,\%$ \\
$M$ & 1 & $1031$ & $1046$ & $5\,\%$ \\
$M$ & top of group & $\approx 2710$ & $2728$ & $95\,\%$ \\
\hline
\end{tabular}
\end{center}

\paragraph{Two points that are easy to get wrong.}
\begin{enumerate}
  \item Band 1 never reaches $f_r$ in this diagram, because the zone ends first.
        It tops out at $1031\,\mathrm{Hz}$ at $M$. The limit $f \to f_r$ exists
        only in the unfolded model; in the diagram it corresponds to band 1
        near $M$, with a plate share of $5\,\%$.
  \item The ``upper band at $\Gamma$'' of the two-level model \emph{is} the flat
        band of the plot. It sits at about $1270\,\mathrm{Hz}$, above the red
        line. Nothing in the diagram sits on $1200\,\mathrm{Hz}$.
\end{enumerate}

\paragraph{Outlook.}
The gap runs from the top of band 1 at $M$ to band 2 at $X$. That upper edge is
the $100\,\%$ plate mode, which does not feel the resonator at all. The next
chapter takes this up.

% ---------------------------------------------------------------------
\section{Check questions with answers}

\paragraph{Question 1.} \emph{Just above $f_r$, $m_\mathrm{eff}$ is large and
negative. What does the plate feel there, and why does that rule out a
propagating bending wave?}

\emph{Answer.} The resonator oscillates opposite to the plate and pulls it back,
so the usual propagating wave is stopped. Two refinements:
\begin{itemize}
  \item Nothing is dissipated, because the model is undamped. The wave becomes
        \emph{evanescent}: $\kappa^4 < 0$ makes $\kappa$ complex, so the plate
        still moves, but the motion decays from cell to cell. The energy is
        reflected, not absorbed.
  \item Negative $m_\mathrm{eff}$ alone is not enough. $m_\mathrm{eff}$ is
        negative at \emph{every} frequency above $f_r$, but the gap ends at
        $f_r\sqrt{1+\mu}$. What counts is the \emph{total} mass per area,
        $\rho h + m_\mathrm{eff}/A$. It is negative only while the negative
        mass of the resonator outweighs the mass of the plate, and it crosses
        zero exactly at $\omega^2 = \omega_r^2(1+\mu)$.
\end{itemize}

\paragraph{Question 2.} \emph{Why is the resonator band flat, and where would you
expect it not to be flat?}

\emph{Answer.} The plate moves little there, so the frequency barely depends on
$\mathbf{k}$. What decides flatness is not the distance from $\omega_r$ itself,
but how far the \emph{plate band at that $\mathbf{k}$} is from the resonator
level: slope of a band = plate share $\times$ slope of the plate curve. The
band bends at the avoided crossing, and it moves away from $\omega_r$
\emph{because} it hybridizes. The band near $\Gamma$ shows the difference: it
sits at a frequency different from $f_r$ ($1342\,\mathrm{Hz}$ in the model,
about $1270\,\mathrm{Hz}$ in FE) and is still nearly flat, because the plate
band is far below it there.

\paragraph{Question 3.} \emph{Along the lower band, the eigenvector goes from
mostly plate to mostly resonator. What does this mean for MAC tracking, and
which branch does a shaker on the finite plate see?}

\emph{Answer.} The 50/50 mix sits at the crossing, not near $f_r$; the mode
becomes resonator-dominated towards the end of band 1.

\emph{MAC tracking:} band 1 goes from $80\,\%$ plate at $\Gamma$ to $5\,\%$ at
$M$. That is a large change of character, but it happens gradually over many
$\mathbf{k}$-points, so the MAC between neighbouring points stays high and
tracking should follow the band without swapping labels. The MAC must use the
\emph{full vector including the resonator DOF}, with the augmented mass matrix;
on plate DOFs alone the comparison is somewhat weaker (lowest MAC of band 1:
$0.969$ instead of $0.990$), and leaving out the mass matrix is far worse ($0.619$).

\emph{Shaker:} the question has no clean one-branch answer. The shaker drives the
plate, so a mode is excited in proportion to its plate share. It sees band 1 well
below about $850\,\mathrm{Hz}$, band 1 near $M$ hardly at all, and the flat part
of band 2 only weakly (reasoning from the model, not measured). Near the gap
edges the upper branch keeps more plate motion than the lower one: at the upper
edge at $\Gamma$ the model gives a plate share of $\mu/(1+\mu) = 20\,\%$ and
$u/w = -1/\mu = -4$.
